\documentclass[11pt,a4paper]{article}
\usepackage[utf8x]{inputenc}
\usepackage[T1]{fontenc}
\usepackage[spanish]{babel}
\usepackage{amsmath}
\usepackage{amssymb,amsfonts,textcomp}
\usepackage{xcolor}
\usepackage{array}
\usepackage{multirow}
\usepackage{hhline}
\PassOptionsToPackage{hyphens}{url}\usepackage{hyperref}
\usepackage{float}
\usepackage{xkeyval}
\usepackage[pdftex]{graphicx}
\usepackage[yyyymmdd,hhmmss]{datetime}
\usepackage{appendix}
\usepackage{listings}
\usepackage{wasysym}
%\usepackage[centering,width=8.5in,height=11in]{geometry}
\usepackage[centering,a4paper]{geometry}
\usepackage{pstricks}
\usepackage{tikz}
\usetikzlibrary{calc, backgrounds}
\usepackage{pdfpages}
\definecolor{dkgreen}{rgb}{0,0.6,0}
\definecolor{gray}{rgb}{0.5,0.5,0.5}
\definecolor{mauve}{rgb}{0.58,0,0.82}
\definecolor{lstbackground}{rgb}{0.95,0.95,0.95}
%\lstset{frame=tb,
%	backgroundcolor=\color{lstbackground},
%  language=Bash,
%  aboveskip=3mm,
%  belowskip=3mm,
%  showstringspaces=false,
%  columns=flexible,
%  basicstyle={\small\ttfamily},
%  numbers=none,
%  numberstyle=\tiny\color{gray},
%  keywordstyle=\color{blue},
%  commentstyle=\color{dkgreen},
%  stringstyle=\color{mauve},
%  breaklines=true,
%  breakatwhitespace=true
%  tabsize=3
%}

\lstset{frame=tb,
	backgroundcolor=\color{lstbackground},
%  language=Bash,
  aboveskip=3mm,
  belowskip=3mm,
  showstringspaces=false,
  columns=flexible,
  basicstyle={\small\ttfamily},
  numbers=none,
%  numberstyle=\tiny\color{gray},
%  keywordstyle=\color{blue},
%  commentstyle=\color{dkgreen},
%  stringstyle=\color{mauve},
  breaklines=true,
  breakatwhitespace=true
  tabsize=4
}

\usepackage{verbatim}
\begin{comment}
	\hypersetup{
		pdftex, 
		colorlinks=true, 
		linkcolor=blue, 
		citecolor=blue, 
		filecolor=blue, 
		urlcolor=blue, 
	pdftitle={Software Libre}, 
	pdfauthor={Eduardo Grosclaude}, 
	pdfsubject={Documento de la materia Software Libre}, 
	pdfkeywords={Software Libre, Tecnicatura en Administración de Sistemas y Software Libre, Universidad Nacional del Comahue}
	}
\end{comment}	

%\addto\captionsspanish {%
%	\def\appendixname{Apéndices}
%}
% Outline numbering
\setcounter{secnumdepth}{1}
% Reset section numbering between parts
\makeatletter
\@addtoreset{section}{part}
\makeatother  
% List styles
\newcommand\liststyleLi{%
\renewcommand\labelitemi{\tiny${\blacksquare}$}
\renewcommand\labelitemii{\tiny${\square}$}
\renewcommand\labelitemiii{\tiny${\circ}$}
\renewcommand\labelitemiv{\tiny${\circ}$}
}
\newcommand\liststyleLii{%
\renewcommand\labelitemi{{\textbullet}}
\renewcommand\labelitemii{${\circ}$}
\renewcommand\labelitemiii{${\blacksquare}$}
\renewcommand\labelitemiv{{\textbullet}}
}
\newcommand\liststyleLiii{%
\renewcommand\labelitemi{{\textbullet}}
\renewcommand\labelitemii{${\circ}$}
\renewcommand\labelitemiii{${\blacksquare}$}
\renewcommand\labelitemiv{{\textbullet}}
}

\liststyleLi

% Page layout (geometry)
\setlength\voffset{-1in}
\setlength\hoffset{-1in}
\setlength\topmargin{2cm}
\setlength\oddsidemargin{2cm}
\setlength\textheight{23.246668cm}
\setlength\textwidth{17.006cm}
\setlength\footskip{26.144882pt}
\setlength\headheight{1.016cm}
\setlength\headsep{0.508cm}
% Footnote rule
\setlength{\skip\footins}{0.119cm}
\renewcommand\footnoterule{\vspace*{-0.018cm}\setlength\leftskip{0pt}\setlength\rightskip{0pt plus 1fil}\noindent\textcolor{black}{\rule{0.25\columnwidth}{0.018cm}}\vspace*{0.101cm}}
% Pages styles
\makeatletter
\newcommand\ps@Standard{
  \renewcommand\@oddhead{{\raggedleft Cabecera \ } {\raggedright \thepage{}}}
  \renewcommand\@evenhead{\@oddhead}
  \renewcommand\@oddfoot{}
  \renewcommand\@evenfoot{\@oddfoot}
  \renewcommand\thepage{\arabic{page}}
}
\makeatother
% \pagestyle{Standard}
\usepackage{fancyhdr}
%\usepackage{libertine}

\usepackage{lmodern}
\renewcommand*\familydefault{\sfdefault} %% Only if the base font of the document is to be sans serif


% \renewcommand*\familydefault{\sfdefault}
\pagestyle{fancy}
% footnotes configuration
\makeatletter
\renewcommand\thefootnote{\arabic{footnote}}
\makeatother
\title{Redes II}
\author{Eduardo Grosclaude}
\date{2014-08-13}
\usepackage{graphicx}

\usepackage{xkeyval}
\usepackage{pifont}
\usepackage{xcolor}
\newcommand{\revisar}[1]{{\color{red}[#1]}}



%\usepackage[hundred]{vrsion}

\newcommand{\borrador}{
\revisar{\today, \currenttime  -  Material en preparación, se ruega no imprimir mientras aparezca esta nota}
}


%\newcommand{\nota}[1]{{\color{red}[#1]}}
%\newcommand{\revisar}[1]{}
\newcommand{\nota}[1]{}

\newcommand{\nonota}[1]{#1}

\newcommand{\quotes}[1]{``#1''}

   
\newcommand{\shade}[1]{\textcolor{black!50}{#1}}

% ancho opcional, por defecto 15cm
% \figura{copyleft}{Símbolo de Copyleft}{copyleft.png}
% \figura[6]{copyleft}{Símbolo de Copyleft}{copyleft.png}
\newcommand{\figura}[4][15]{
 \begin{figure}[htbp] 
 \centering 
 \includegraphics[width=#1cm]{./img/#4} 
 \caption{#3} 
 \label{fig:#2} 
 \end{figure} 
}

\newcommand{\tabla}[4]{
 \begin{table} 
 \centering 
 \small
 \begin{tabular}{#3}
 #4
 \end{tabular}
 \caption{#2}
 \label{tab:#1} 
 \end{table} 
}

\usepackage{mdframed}

\newcommand{\recuadro}[1]{\vspace*{0.2cm}
\begin{minipage}[l]{0.84\textwidth}
\begin{mdframed}
#1
\end{mdframed}
\end{minipage}\vspace*{0.4cm}
}

% \newcommand{\recuadro2}[1]{
% \fcolorbox{black}{grey}{
% \parbox[t]{1.0\linewidth}{ \vspace*{0.4cm}#1\vspace*{0.4cm} } 
% }
% 
\usepackage{sectsty}
%\usepackage[compact]{titlesec} 
\definecolor{bl}{rgb}{0.0,0.2,0.6} 
\allsectionsfont{\color{bl}\scshape\selectfont}


\newcommand{\code}[1]{\lstinline$#1$}
%\newcommand{\code}[1]{\begin{verbatim}{#1}\end{verbatim}}

\hypersetup{colorlinks=true, linkcolor=blue, citecolor=blue, filecolor=blue, urlcolor=blue, 
	pdftitle={Redes II}, 
	pdfauthor={Eduardo Grosclaude}, 
	pdfsubject={Documento de la materia Redes II}, 
	pdfkeywords={Redes, Switching, Alta Disponibilidad, Tecnicatura en Administración de Sistemas y Software Libre, Universidad Nacional del Comahue}}

%TODO Incorporar load balance en LAN con modulo nth de iptables

% --------------------------------------------------------------------
\begin{document}\sloppy
%% \pagestyle{empty}\psset{unit=1in}\begin{pspicture}(8.27,11.69)% use your page size
%   \rput[b](3.5,8){\parbox{5in}{\begin{flushright}\Huge\bfseries\sffamily\color{blue}{Redes II TUASSL}\end{flushright}}}
%   \uput[-90](3.5,8){\color{blue}\rule{5in}{1ex}}
% 	\noindent\includegraphics[width=\paperwidth]{img/Patrik_Goethe_Unsplash.jpg}
% \end{pspicture}
% 



% 
%\usepackage{tikz}
%\usetikzlibrary{calc}
%
% argumentos radio, x,y, transparencia
%
\newcommand{\logofai}[4] {
\begin{tikzpicture} %[remember picture,overlay]
\begin{scope}
 	\def \rmenor {#1/3*2cm}
	\def \rmayor {#1cm}
	\def \x {#2cm}
	\def \y {#3cm}
	
	\fill[blue,color=blue] (\x,\y) circle(\rmayor);
	\fill[white,color=white] (\x,\y+\rmayor) rectangle(\x+\rmayor,\y-\rmayor);
	\fill[white,color=white] (\x,\y) circle(\rmenor);
	\fill[blue,color=blue] (\x-\rmayor/10,\y+\rmayor/10) circle (0.44*\rmayor);

	\def \gap {35}
	\def \aini {90}
	\def \afin {115}
  	\foreach \s in {1,...,5} {

	  	\draw[fill=gray, color=gray!80]
	  	([shift={(\aini-\s*\gap:\rmenor)}] \x,\y) 
	  		arc(\aini-\s*\gap:\afin-\s*\gap:\rmenor) -- 
  		([shift={(\afin-\s*\gap:\rmayor)}] \x,\y) 
  			arc(\afin-\s*\gap:\aini-\s*\gap:\rmayor) -- 
  		cycle;

	}
	\draw[fill=white,color=white]
  		([shift={(20-3*\gap:\rmenor)}] \x,\y) arc(20-3*\gap:45-4*\gap:\rmenor)-- 
  		([shift={(45-4*\gap:\rmayor)}] \x,\y) arc(45-4*\gap:20-3*\gap:\rmayor)-- 
  		cycle;
  		
  	%% Transparencia
  	\def \size {#1 * 2cm + 0.05cm}
  		\begin{tikzpicture}[remember picture,overlay]
			\node[fill=white,color=white,opacity=#4,
							minimum size=\size,circle] at (\x,\y) {};
		\end{tikzpicture}

	\clip(\x,\y) circle(\rmayor + 0.05cm);
	\end{scope}
\end{tikzpicture}
}

% \begin{document}
% 
%     % 4 cm radio, ubicado en (3,9), 0.3 de transparencia
% 	\logofai{4}{3}{9}{0.3}
% 	\logofai{3}{1}{-2}{0}
% 	\logofai{1}{-3}{-3}{0.5}
% \end{document}

% \logofai{4}{3}{9}{0.3}
% 	%\includegraphics[trim = 10mm 80mm 20mm 5mm, clip, width=3cm]{chick}
% 	%\includegraphics[width=100]{img/Patrik_Goethe_Unsplash.jpg}
% 
% 
% %\usepackage{pdfpages}
% \pagestyle{empty}
% \includegraphics[width=\paperwidth]{img/RII-portada.pdf}
% 
% \begin{minipage}
% {\paperwidth}	
% \includepdf{img/RII-portada.pdf}
% \end{minipage}
% 
% Found the solution here: http://tex.stackexchange.com/questions/5911/how-to-include-pdf-pages-without-a-newpage-before-the-first-page
% 
% the only thing is to add the \section in the pagecommand of \includepdf like this:
% 
% \includepdf[pages=-,scale=.8,pagecommand={\section{title 1}\label{}},linktodoc=true]{myPDF1.pdf}
% \includepdf[pages=-,scale=.8,pagecommand={\section{title 2}\label{}},linktodoc=true]{myPDF2.pdf}

% %\begin{tikzpicture}[remember picture,overlay]
%  \begin{tikzpicture}
%    \node at (current page.center) {\includegraphics[page=1]{img/RII-portada.pdf}};
%  \end{tikzpicture}

% \includepdf[pages=2-last]{MYINCLDOC.pdf} 
% Nada anduvo :(


 :(
\maketitle

\borrador

\abstract {En este escrito se presenta la descripción y material inicial de la asignatura \textbf{Redes II}, para la carrera de Tecnicatura Universitaria en Administración de Sistemas y Software Libre, de la Universidad Nacional del Comahue. 

La materia es electiva, cuatrimestral en modalidad presencial y las clases son de carácter teórico-práctico, desarrolladas en forma colaborativa. Está preparada con los objetivos generales de  \textbf{actualizar y ampliar el instrumental de trabajo en ambientes de redes}. 
  


\newpage
\emph{" "}
\newpage

\tableofcontents

\newpage
\emph{" "}

%----------- P R E S E N T A C I O N  ---------
\newpage
\part {La asignatura}


\section{Objetivos}
\subsection{De la carrera}
Según el documento fundamental de la Tecnicatura, el Técnico Superior en Administración de Sistemas y Software Libre estará capacitado para:
\begin{itemize}
	\item Desarrollar actividades de administración de infraestructura. Comprendiendo la administración de sistemas, redes y los distintos componentes que forman la
infraestructura de tecnología de una institución, ya sea pública o privada.
	\item Aportar criterios básicos para la toma de decisiones relativas a la adopción de nuevas tecnologías libres.
	\item Desempeñarse como soporte técnico, solucionando problemas afines por medio de la comunicación con comunidades de Software Libre, empresas y desarrolladores de
software.
	\item Realizar tareas de trabajo en modo colaborativo, intrínseco al uso de tecnologías libres.
	\item Comprender y adoptar el estado del arte local, nacional y regional en lo referente a implementación de tecnologías libres. Tanto en los aspectos técnicos como legales.
\end{itemize}
\subsection{De la asignatura}

\begin{itemize}
	\item Actualizar y ampliar el instrumental de trabajo en ambientes de redes.
\end{itemize}


\section{Cursado}
\begin{itemize}
	\item Cuatrimestral de 16 semanas, 4 horas semanales, 64 horas totales
	\item Clases teórico-prácticas presenciales
	\item Promocionable con trabajos prácticos
\end{itemize}


\section {Contenidos}
\subsection{Contenidos mínimos}
\begin{itemize} 
	\item Switching. 
	\item Redes Privadas Virtuales. 
	\item Balance de carga y Alta Disponibilidad en redes. 
\end{itemize}


\subsection {Programa}
\begin{enumerate}
	\item Switching
		\begin{itemize}
			\item Red Ethernet, dominio de colisión y dominio de Broadcast
			\item Bridges, switches, arquitectura y funcionamiento. 
			\item VLANs, trunking
		\end{itemize}
	\item Redes Privadas Virtuales
		\begin{itemize}
			\item Encapsulamiento
			\item VPN de nivel 2 y de nivel 3
			\item Configuración y administración de OpenVPN
			\item Configuración de ruteo y de firewalling en VPN
		\end{itemize}
	\item Alta Disponibilidad y Balance de carga en Redes
		\begin{itemize}
			\item Arquitecturas redundantes 
			\item Bonding y modos de configuración
			\item Protocolo STP 802.1d
			\item Balance de carga, Ruteo por políticas
%			\item Clustering de LB 
		\end{itemize}

\end{enumerate}

\section {Bibliografía inicial}
\begin{itemize}
\item C., Zimmerman, Joann Spurgeon, Ethernet switches. Sebastopol, CA: O’Reilly Media, 2013.
\item M. Feilner, A. Bämler, N. Graf, and M. Heller, Open VPN building and integrating virtual private networks. Birmingham, U.K.: Packt Pub., 2006.
\item T. Bourke, Server load balancing, 1st ed. Beijing; Sebastopol, Calif: O’Reilly, 2001.
\end{itemize}

%
\section{Evaluación}
La evaluación de la materia se realizará mediante trabajos grupales de investigación y desarrollo sobre proyectos de Software Libre, de la siguiente manera.
\begin{itemize}
	\item Los estudiantes se dividirán en grupos de 2 a 5 personas. 
	\item Los grupos desarrollarán trabajos prácticos en etapas que se distribuirán a lo largo de la materia. 
	\item Cada grupo abrirá un diario, blog o wiki de acceso público en cualquier sitio disponible y publicará, mediante el Foro de la materia, la forma de acceder al diario para lectura. Los docentes y los demás estudiantes de la materia podrán acceder al diario del grupo para lectura. Todo cambio en la dirección o forma de acceso deberá ser informado mediante el Foro.
	\item El grupo irá aportando los resultados de cada etapa de los trabajos a su diario, y periódicamente comentará además en clase las experiencias surgidas durante la realización de los trabajos.
	\item El material publicado en el diario será reunido en un documento final que será entregado \textbf{en formato electrónico} al finalizar la materia. El documento indicará tema del trabajo, resumen, integrantes del grupo, desarrollo y conclusiones. 
	\item El documento será acompañado por una presentación de no más de treinta minutos que será expuesta según el cronograma adjunto. 
	\item La acreditación final tendrá en cuenta la calidad del material aportado al diario por el grupo, la calidad de los documentos finales de los trabajos, la presentación oral y la participación en clase ofreciendo la experiencia adquirida durante la realización de los trabajos.
\end{itemize}

\subsection {Trabajo I - Colaboración con proyectos libres}
\subsubsection{Etapa 1}  
Descargar e instalar software ofrecido por un proyecto de Software Libre que esté en actividad (puede tratarse de un entorno de escritorio, un programa de sistema, programas de usuario final, una distribución completa, etc.). Familiarizarse con el software utilizándolo. 
\subsubsection{Etapa 2} 
Basándose en el conocimiento adquirido con el uso del software, colaborar de alguna forma con el proyecto que lo origina: 
\begin{itemize}
	\item traduciendo o localizando parte del software,
	\item generando documentación faltante, 
	\item traduciendo parte de la documentación, 
	\item detectando y denunciando errores en el software o en la documentación,
	\item aportando, modificando o corrigiendo código,
	\item aportando conocimiento a los usuarios del proyecto en blogs, salas de chat, bases de conocimiento, etc.
\end{itemize}
Puede abordarse cualquier cantidad manejable de proyectos. La colaboración debe consistir en alguna interacción positiva y completa con cada proyecto. El grupo incorporará al diario los reportes que acrediten esa interacción. Cuando no sea posible realizar o completar la interacción se indicarán las causas, y las acciones realizadas.

El aporte al proyecto debe efectuarse por los canales establecidos por el proyecto. Si se trata de documentación, respetar el formato utilizado; si es el reporte de un error, hacerlo por la vía preferida por el proyecto, etc.

\subsubsection{Etapa 3} 
El grupo entregará un documento conteniendo la historia de las interacciones con cada proyecto, adjuntando las pruebas en anexos y ofrecerá una presentación.

\subsection {Trabajo II - Evaluación de proyectos libres}

\subsubsection{Etapa 1} 
El grupo enunciará un determinado requerimiento concreto de software que puede ser presentado por un empleador. Algunos ejemplos posibles son:
\begin{itemize}
	\item \quotes{un servidor de correo electrónico que maneje listas},
	\item  \quotes{una aplicación de control de asistencia para empleados},
	\item  \quotes{un sistema de edición de textos para traductores},
	\item  \quotes{un sistema de gestión de contenidos web que incluya workflow}, 
	\item \quotes{un motor de juegos 2D para crear juegos que asistan en la enseñanza de matemática},
	\item  \quotes{un programa de simulación de ataques para evaluar postura de seguridad}, 
	\item \quotes{un sistema de control de stock para zapaterías},
	\item \quotes{una distribución de GNU/Linux para escuelas de arte},
	\item \quotes{una distribución para sistemas empotrados}, etc.
\end{itemize}
El grupo debe comprender el propósito del software requerido y debe contar con al menos un integrante con conocimiento razonable de la temática involucrada. El grupo escribirá una entrada en el diario consignando toda la información posible sobre los requerimientos. 

\subsubsection{Etapa 2} 

\begin{itemize}
	\item El grupo $n$ (en adelante \quotes{el proveedor}) tomará a su cargo el requerimiento del grupo $n+1$ (en adelante \quotes{el cliente}), y se atendrá a dicha descripción para el resto del trabajo. 
	\item El grupo proveedor buscará proyectos de SL que apunten a cubrir esos requerimientos y seleccionará al menos dos proyectos, idealmente tres, de entre ellos.
\end{itemize}

\subsubsection{Etapa 3}
Los proyectos serán comparados en función de varios parámetros o dimensiones.
\begin{itemize}
	\item  ajuste a los requerimientos (actual, previsto o potencial),
	\item  licenciamiento, 
	\item  motivación del desarrollo, 
	\item  modelos de negocio del proyecto, 
	\item  tamaño y permanencia de la comunidad,
	\item  dinámica de soporte, 
	\item  dinámica de actualizaciones y mejoras del software.
\end{itemize}

Se podrán agregar a la comparación uno o más desarrollos no libres. 

Las dudas sobre detalles de los requerimientos serán dirigidas al grupo cliente, y contestadas por aquél, mediante el Foro de la página de la materia.  
\subsubsection{Etapa 4} 
El grupo entregará un documento conteniendo la comparación y haciendo una recomendación final, explicando sus fundamentos. Deberán volcar en el trabajo lo que se vaya aprendiendo durante el curso de la materia, en cada uno de los parámetros o dimensiones nombrados. Finalmente ofrecerán una presentación sobre el trabajo.

\label{sub:acreditacion}

\section {Cronograma}
\begin{tabular}{c|c|l|l|l}
Fecha & Semana & Unidad & Trabajo I & Trabajo II\\
\hline
\hline
14/8 & 1	& 	1. Switching 			& 	& \\
21/8 & 2 	& 				 	& 	& \\
28/8 & 3	& 					&	& \\
4/9 & 4 	& 					&	& \\
\hline
\hline
11/9 & 5	& 	2. Redes Privadas Virtuales 	& 	& \\
18/9 & 6	& 					&	& \\
25/9 & 7 	& 					&   & \\
2/10 & 8	& 					&  	& \\ 
9/10& 9		& 	Parcial I			& 	& \\
\hline
\hline
16/10 & 10	& 	3. Balance de Carga en Redes	& 	& \\ 
23/10 & 11	& 					& 	& \\
30/10 & 12	& 					&	& \\
6/11 & 13	& 					&	& \\
\hline
\hline
13/11 & 14	& 	4. Alta Disponibilidad en Redes	&	& \\
20/11 & 15	& 					&	& \\
27/11 & 16	& 	Parcial II			&	& \\ 
\hline
\end{tabular}



% \begin{tabular}{|r|c|c|c|c|c|c|c|c|}
% \hline
%\textsf{7} & fbox {algo} & & & & & & &\\ 
%\hline
%\textsf{7} & & & & & & & &\\ 
%\hline
%\end{tabular}

% subsection  (end)


%----------- M A T E R I A L ---------
\newpage
\part {Switching}



\section{Ethernet}
\begin{enumerate}
\item ¿Qué módulos del kernel Linux están controlando las interfaces de red de su equipo? ¿Cuáles comandos permiten conocer las marcas y modelos de las placas instaladas? ¿Qué diferencia hay entre placas y chipsets? ¿Cómo se puede investigar cuál módulo corresponde a cuál marca y modelo de placa de red? 
\item ¿Qué es el bus PCI y qué son los números de dispositivos asignados? ¿Cuál es la relación entre el bus PCI y los nombres de dispositivos de red en Linux?
\item ¿Qué información aporta el comando ethtool? ¿Qué clases de placas de red soporta? ¿Qué significa \emph{Link detected} en la salida de este comando?
\item ¿Qué información transporta un frame Ethernet? ¿Qué diferencias existen entre los frames que circulan por una Ethernet cableada y una inalámbrica? 
\item ¿Qué diferencias existen entre el diseño de la red Ethernet original, implementada con coaxil, y la que está implementada con cableado en estrella, conectada mediante hubs? ¿Qué diferencias existen entre esta última y una implementada con switches? 
\end{enumerate}


\section{Dominios de colisión y de broadcast}
\label{sub:}
\begin{enumerate}
\item ¿Qué es un bridge, qué configuración lleva, qué función realiza sobre los frames que circulan por la red, y cómo aprende su información de trabajo? 
\item ¿Existen bridges inalámbricos? 
\item ¿A qué se llama segmentación?
\item ¿Un hub es equivalente a un bridge? ¿Un router es equivalente a un bridge? ¿Un switch es equivalente a un bridge? 
\item ¿Qué es un frame de broadcast? ¿Qué es un frame unicast? ¿Cómo se distingue un frame de broadcast? 
\item ¿Qué protocolos utilizan frames de broadcast? ¿Cómo es aprovechado el mecanismo de broadcast por el protocolo ARP? ¿Qué otros usos podría recibir el mecanismo de broadcast?
\item ¿Qué conducta observan los hosts de la red al recibir un frame de broadcast? ¿En qué casos puede resultar esto un problema y por qué? ¿Cuáles son las contramedidas?
\item ¿Qué es un dominio de colisión?  ¿Qué es un dominio de broadcast? ¿Qué tipo de dispositivo es el que delimita la frontera entre dos dominios de colisión?  ¿Qué tipo de dispositivo es el que delimita la frontera entre dos dominios de broadcast? 
\item ¿Qué efecto tiene la introducción de un bridge en un dominio de colisión? El número de equipos en un mismo dominio de colisión ¿se incrementa o se reduce? ¿Qué efecto tiene esto sobre la probabilidad de colisión?
\item ¿Qué diferencia hay entre un paquete de broadcast y un frame de broadcast? Los routers, ¿repiten los frames de broadcasts?
\item En la figura \ref{fig:broadcast}, ¿cuáles son los dominios de colisión y de broadcast que pueden distinguirse? Si el host H1 emite un frame de broadcast, ¿qué otros hosts lo reciben? Si el host H4 emite un broadcast, ¿qué otros hosts lo reciben?
\item En la figura \ref{fig:arbol}, ¿cuáles son los dominios de colisión y de broadcast que pueden distinguirse? Si el host H1 emite un frame de broadcast, ¿qué otros hosts lo reciben? 
\item El concepto de dominio de colisión, ¿sigue siendo aplicable en la actualidad, cuando prácticamente la totalidad de las redes usan bridging en lugar de repetidores? 
\item El uso de bridges, ¿limita el tráfico de broadcast sobre un segmento de red?
 \end{enumerate}
% subsection  (end) Dominios de colisión y de broadcast


\figura[12]{broadcast}{Colisiones y broadcasts}{broadcast.jpg}


\section{Switches}
\label{sub:Switches}

\begin{enumerate}
\item ¿Cuál es la función de un switch? 
\item ¿Cuántos dominios de colisión hay en una red implementada sobre un switch? ¿Cuántos dominios de broadcast hay en una red implementada sobre un switch?
\item ¿En qué consiste el proceso de conmutación de un frame? ¿Qué modos de conmutación son los más usados?
\item ¿Cómo aprende su información de trabajo un switch? ¿Dónde almacena dicha información? ¿Cómo accede a esta información?
\item ¿Qué debe hacer un switch ante un frame dirigido a una estación (\quotes{unicast}) cuya dirección destino figura en sus tablas? 
\item ¿Qué debe hacer un switch ante un frame dirigido a una estación (\quotes{unicast}) que no figura en ninguna de sus tablas? ¿Cómo se llama esta acción?
\item ¿Qué debe hacer un switch ante un frame cuya dirección origen no figura en sus tablas? 
\item ¿Qué debe hacer un switch ante un frame de broadcast? ¿Cómo se llama esta acción?
\item En una topología compuesta por un árbol de tres switches como en la figura \ref{fig:arbol}, ¿cómo aprenden los switches las direcciones de los equipos que no están directamente conectados? ¿Cuántos dominios de broadcast aparecen en dicha figura?
\item ¿En qué orden aparecen las direcciones destino y origen en el frame Ethernet? ¿En qué orden aparecen las direcciones destino y origen en los paquetes IP y en los segmentos TCP o UDP? ¿De qué forma aprovechan los switches la diferencia?

\item En la figura \ref{fig:dominios}, si el host H1 establece conexión TCP con H3, ¿quiénes más pueden ver el tráfico entre H1 y H3? Si el host H1 establece conexión TCP con H4, ¿quiénes más pueden ver el tráfico entre H1 y H4?

\item ¿Cuántas redes aparecen, respectivamente, en las topologías de figuras \ref{fig:broadcast}, \ref{fig:arbol} y \ref{fig:dominios}? ¿Cuántos espacios de direccionamiento IP diferentes habrá en cada topología?

\item ¿Existen switches inalámbricos? ¿En qué se parecen y en qué se diferencian un switch cableado y un punto de acceso inalámbrico? 
\item ¿Qué diferencias existen entre las conexiones de switches en cascada y apiladas (\emph{stacked})? 
\item ¿Qué velocidad suelen tener los ports de un switch? ¿Qué dicen las normas de cableado respecto de la calidad del cable, su longitud y su forma de instalación, para cada caso?
\item Si los ports de un switch tienen una velocidad determinada, ¿tiene sentido conectar a un port un conjunto de hosts capaces de transmitir, en forma agregada, mayor cantidad de bits por segundo que la velocidad admitida por el port? ¿Qué significa \emph{oversubscription}?
\item En una red compartida implementada con hubs, ¿cuál es el ancho de banda obtenido por cada cliente? ¿Y en una red implementada con switches?

\item En la Fig. \ref{fig:velocidades}, supongamos que los hosts H1 y H3 mantienen tráfico entre sí al mismo tiempo que H2 y H4 lo hacen entre sí. ¿Cuál será la velocidad máxima de dicho tráfico si el dispositivo es un hub que funciona a 100Mbps? ¿Y si es un switch que funciona a 100Mbps? 

\item En aquellos switches con ports de diferentes velocidades, ¿cuáles podrían ser los usos de los ports de mayor velocidad?
\item ¿Qué son los dispositivos de red administrables? ¿En qué consiste el protocolo SNMP? ¿Qué características operativas de un switch administrable pueden consultarse o modificarse mediante SNMP? 
\item ¿Qué características administrables de un switch tienen que ver con la gestión de seguridad en redes?
\end{enumerate}
% subsection  (end) Switches
\figura[8]{velocidades}{Competencia por el medio}{velocidades.jpg}



\figura[10]{arbol}{Jerarquía de switches}{arbol.jpg}

\figura[12]{dominios}{Dominios}{dominios.jpg}


\section{Arquitecturas de redundancia}
\begin{enumerate}
\item ¿Qué ventajas aporta la redundancia de enlaces en una red LAN? ¿Qué formas de redundancia se pueden implementar? ¿En qué consiste la tolerancia a fallos?
\item ¿Qué ocurre en una red Ethernet con múltiples caminos entre dos hosts? ¿Qué escenarios de fallo se pueden presentar?  
\item ¿Cómo se puede detectar una tormenta de broadcasts?
\item ¿Qué ventaja ofrece el protocolo IEEE 802.1d (STP)? ¿Cómo es su modo de operación y qué intenta construir? ¿Qué dispositivos soportan 802.1d? 
\item ¿Cómo se elige un dispositivo raíz del árbol STP? ¿Cómo se determinan los caminos al raíz? ¿Qué demora puede tener normalmente la convergencia de STP?
\item ¿Qué datos de configuración o estado, que sean relevantes para 802.1d, y cuya configuración por el administrador sea posible, admiten los bridges y switches? ¿En qué puede influir el administrador de redes sobre la elección del raíz? ¿Cómo se puede aprovechar este hecho?
\item ¿En qué consiste el protocolo RSTP?
\end{enumerate}
% subsubsection  (end) Arquitecturas de redundancia



\section{VLANs}
\label{sub:}
Un diseño clásico de redes de campus consiste en un router que proyecta segmentos de LAN como radios de una estrella (Fig. \ref{fig:bbcolap}). La función de comunicar los radios, que anteriormente era cumplida por el backbone de la red, en este diseño se logra por la conmutación efectuada por el router, por lo cual suele llamarse de backbone colapsado. 

Con este diseño, los dominios de broadcast, y por lo tanto los espacios IP definidos sobre ellos, quedan limitados geográficamente. Si en una zona del campus donde llega un radio de la estrella se necesita ubicar nodos sobre dos dominios de broadcast (porque se desea aislarlos por motivos de seguridad, porque se desea situar equipos sobre dos espacios IP diferentes, o porque se desea limitar la competencia de ambas clases de tráfico por el medio), debe haber dos cableados y deben ocuparse dos bocas del router central. 

\figura[8]{bbcolap}{Backbone colapsado}{backbone-colapsado.jpg}

\figura[12]{vlans}{VLANs}{vlans.jpg}


Con la funcionalidad avanzada de VLANs provista por algunos switches (y definida en el estándar IEEE 802.1Q), el mismo cableado, y el sistema de switches de llegada, puede usarse para conducir dos o más dominios de broadcast. 
 
\begin{itemize}


	\item ¿En qué consiste el diseño de red de campus de backbone colapsado? ¿Qué impacto tiene este diseño sobre la posibilidad de distribuir equipos de diferentes clases sobre una misma región de la red?
	\item 
¿Cuál es la finalidad de definir VLANs en un switch? 

	\item 
¿Cómo se modifica el formato de frame Ethernet para lograr la capacidad de separar los dominios de broadcast al definir VLANs?

	\item ¿Qué debe hacer un switch con VLANs definidas al recibir un frame de broadcast sobre una de sus interfaces? ¿Qué debe hacer con un frame unicast?

	\item En la topología de la figura \ref{fig:vlans}, los tres switches tienen definidas dos VLANs. Los hosts H1, H2 y H4 pertenecen a la VLAN 1, y los hosts H3, H5 y H6 a la VLAN 2. 

	\begin{itemize}
		\item ¿Qué deben hacer los switches con un frame de broadcast recibido desde el host H2?
		\item ¿Qué deben hacer los switches con un frame unicast recibido desde el host H5 y dirigido a H6? ¿Lo mismo, pero desde H5 a H3? ¿Qué diferencia en el formato de los frames existe entre un caso y otro, en cada punto del camino?
		\item ¿Qué condición deben cumplir los ports que interconectan los switches entre sí para poder distribuir las VLANs por toda la topología?
		\item ¿Es posible que una aplicación en el host H5 inicie conexión TCP/IP con una aplicación servidora situada en H2? 
	\end{itemize} 

\item ¿Qué elemento externo es necesario para conectar diferentes VLANs en una misma jerarquía de switches? 
\item ¿Qué son los switches multicapa o \emph{multilayer}? 
\end{itemize} 

% subsection  (end) VLANs


\section {Estudio de caso I}


Una organización desarrolla sus actividades en un campus con dos edificios principales, que alojan las oficinas y talleres. En la organización existen tres áreas principales: Operaciones, Comercialización, e Ingeniería. 

Cada área utiliza un servidor de base de datos principal que es propio del área. La organización cuenta además con un servidor de web y de correo electrónico, un servidor de archivos y un servidor de backups, los tres de uso general para las tres áreas. Se desea que todos los puestos de trabajo puedan acceder además a Internet. 

Una preocupación especial de la organización es darle protección a los servidores y puestos de trabajo de Ingeniería, que no deben ser accedidos desde las demás áreas.
 
La planta del campus y sus edificios, con los principales puntos donde se necesita conectividad, es como indica la Fig. \ref{fig:caso01}. En este diagrama se muestran, con diferentes colores, los puestos de trabajo de cada área. 


\figura[12]{caso01}{Conectar el campus de una organización}{caso01.jpg}


En base a esta información, indique:
\begin{itemize}
	\item Dónde situaría los servidores mencionados.
	\item Dónde ubicaría los elementos de conectividad (switches y routers).
	\item Qué cantidad de puertos debería tener cada uno de estos elementos.
	\item Con qué medios (cobre, fibra, inalámbricos) vincularía los clientes y los elementos de conectividad.
	\item Cómo distribuiría direcciones IP para los clientes y servidores. 
	\item Si utilizaría alguna arquitectura de VLANs, cuál, y por qué. En caso afirmativo, cómo relacionaría las VLANs entre sí, qué enlaces entre switches y routers deben ser de trunking y por qué.
	\item Si utilizaría alguna forma de redundancia, y en caso afirmativo, cuál debe ser el camino normal del tráfico y por qué.
\end{itemize}









\section{VLANs}
Un diseño clásico de redes de campus consiste en un router que proyecta segmentos de LAN como radios de una estrella (Fig. \ref{fig:bbcolap}). La función de comunicar los radios, que anteriormente era cumplida por el backbone de la red, en este diseño se logra por la conmutación efectuada por el router, por lo cual suele llamarse de backbone colapsado. 

Con este diseño, los dominios de broadcast, y por lo tanto los espacios IP definidos sobre ellos, quedan limitados geográficamente. Si en una zona del campus donde llega un radio de la estrella se necesita ubicar nodos sobre dos dominios de broadcast (porque se desea aislarlos por motivos de seguridad, porque se desea situar equipos sobre dos espacios IP diferentes, o porque se desea limitar la competencia de ambas clases de tráfico por el medio), debe haber dos cableados y deben ocuparse dos bocas del router central. 

\figura[8]{bbcolap}{Backbone colapsado}{backbone-colapsado.jpg}

\figura[12]{vlans}{VLANs}{vlans.jpg}


Con la funcionalidad avanzada de VLANs provista por algunos switches (y definida en el estándar IEEE 802.1Q), el mismo cableado, y el sistema de switches de llegada, puede usarse para conducir dos o más dominios de broadcast. 
 
\begin{itemize}


	\item ¿En qué consiste el diseño de red de campus de backbone colapsado? ¿Qué impacto tiene este diseño sobre la posibilidad de distribuir equipos de diferentes clases sobre una misma región de la red?
	\item 
¿Cuál es la finalidad de definir VLANs en un switch? 

	\item 
¿Cómo se modifica el formato de frame Ethernet para lograr la capacidad de separar los dominios de broadcast al definir VLANs?

	\item ¿Qué debe hacer un switch con VLANs definidas al recibir un frame de broadcast sobre una de sus interfaces? ¿Qué debe hacer con un frame unicast?

	\item En la topología de la figura \ref{fig:vlans}, los tres switches tienen definidas dos VLANs. Los hosts H1, H2 y H4 pertenecen a la VLAN 1, y los hosts H3, H5 y H6 a la VLAN 2. 

	\begin{itemize}
		\item ¿Qué deben hacer los switches con un frame de broadcast recibido desde el host H2?
		\item ¿Qué deben hacer los switches con un frame unicast recibido desde el host H5 y dirigido a H6? ¿Lo mismo, pero desde H5 a H3? ¿Qué diferencia en el formato de los frames existe entre un caso y otro, en cada punto del camino?
		\item ¿Qué condición deben cumplir los ports que interconectan los switches entre sí para poder distribuir las VLANs por toda la topología?
		\item ¿Es posible que una aplicación en el host H5 inicie conexión TCP/IP con una aplicación servidora situada en H2? 
	\end{itemize} 

\item ¿Qué elemento externo es necesario para conectar diferentes VLANs en una misma jerarquía de switches? 
\item ¿Qué son los switches multicapa o \emph{multilayer}? 
\end{itemize} 

% subsection  (end) VLANs

\section {Redes Inalámbricas}


Las redes inalámbricas, o \textit{wireless}, utilizan interfaces especiales, como las de radio o de rayos infrarrojos, para emitir tramas o frames en forma similar a las interfaces Ethernet que ya conocemos. Las redes wireless basadas en tecnologías de radio comprenden los vínculos satelitales, los vínculos terrestres de microondas, las redes de telefonía celular, y una familia especial de tecnologías de redes locales, a veces llamada \textbf{wireless LAN o WLAN}. Esta última familia de tecnologías, estandarizada por el grupo de trabajo  \textbf{802.11} de IEEE, permite relacionar entre sí, e incorporar a una LAN cableada, \textbf{dispositivos móviles} como notebooks, tablets, smartphones, u otros.

 La transmisión de datos por medios inalámbricos se enfrenta a problemas físicos muy diferentes de los de la transmisión por cables de cobre o de fibra óptica, como los que conforman las redes cableadas. La distancia entre las interfaces de radio, los obstáculos intermedios (que absorben o reflejan la radiación), las fuentes de ruido electromagnético situadas entre las interfaces, las interferencias entre las redes, son todos factores que tienden a limitar rápidamente las prestaciones de las redes WLAN. Las tasas de transferencia efectivas en una red wireless, y las áreas de cobertura, serán, al menos con las tecnologías y las normas WLAN actuales, mucho menores que las de una red Ethernet. Por esto, las normas de ingeniería usualmente aplicadas consideran a las redes WLAN como una forma complementaria de acceso a las redes cableadas, y no como un reemplazo de éstas.

Como consecuencia de estas dificultades físicas, el diseño y la forma de operación de las interfaces WLAN son sumamente complicados. El hecho de que los dispositivos que se desea conectar, generalmente, son móviles, aumenta la complejidad del problema. Si bien la mayor parte de esta complejidad queda oculta al usuario, conocerla con un poco de detalle nos permite mejorar las prestaciones de una red inalámbrica que vayamos a diseñar, instalar o mantener. 

\subsection {Espectro electromagnético}
Las ondas de radio son esencialmente emisiones de energía que viajan por un medio compartido (el espacio). Una onda (en particular, de radio) es un fenómeno físico que queda descripto matemáticamente por su frecuencia y su amplitud. La amplitud de la onda corresponde al nivel de energía transportada por la onda en cada punto de su trayecto, y su frecuencia corresponde a la cantidad de veces por segundo que la amplitud desarrolla su ciclo (partiendo de un nivel central o 0, creciendo hasta un máximo positivo, decreciendo hasta un mínimo, y volviendo al 0). La frecuencia se mide en \textbf{Hz} (equivalente a 1/s).  Por ejemplo, un ciclo que se repite un millón de veces por segundo tiene una frecuencia de 1MHz. Una señal de 2.4GHz tiene un ciclo que se completa 2.400.000.000 veces por segundo. 

Se llama espectro electromagnético al conjunto de todas las posibles radiaciones, ordenadas por su frecuencia. Una señal compuesta por un conjunto de ondas de diferentes frecuencias (una \textbf{banda}) tiene un \textbf{ancho de banda}, que es la diferencia entre la frecuencia más alta y la más baja de las que la integran. Esta noción de ancho de banda es de naturaleza analógica, y, aunque están estrechamente relacionadas, no debe confundirse con la velocidad de transmisión de una interfaz, medida en bits, que suele llamarse \textbf{ancho de banda digital}. 

Si a un receptor llegan en el mismo momento dos o más señales de radio en la misma frecuencia y con una energía aproximadamente similar, la energía aportada por cada una modifica a la otra (ambas emisiones se hacen recíprocamente \textbf{interferencia}). Dependiendo de las condiciones, esto puede hacer que la información que llevan esas señales quede irreconocible para el receptor que debe interpretarla. Por este motivo, para que las comunicaciones sean eficaces, el espectro debe ser administrado como un recurso limitado, asignando frecuencias a usuarios en cada lugar. Cada Estado dispone alguna regulación o legislación para el uso del espectro en su territorio. Por ejemplo, en cada país, ciertas bandas quedan reservadas para servicios públicos esenciales como seguridad o salud; algunas bandas pueden ser usadas por particulares pero requieren licencias pagas, como las usadas por la radio y la televisión abiertas; y otras, están exentas de licencia.  

Diferentes tecnologías utilizan diferentes bandas del espectro, con sus respectivos anchos de banda. Para las tecnologías de WLAN se han utilizado principalmente las bandas, libres de licencia, alrededor de 2.4 GHz y de 5 GHz. Para cada nuevo desarrollo tecnológico aprobado por el grupo de trabajo correspondiente de IEEE, se han definido métodos de codificación, velocidades de transmisión máximas, alcance geográfico, etc. que son técnica y económicamente posibles dentro de esas bandas.  Así se ha llegado a la definición de una familia de estándares como indica el Cuadro \ref{tab:ieee80211}.
 
 \tabla{ieee80211}{Estándares 802.11 y sus características}{|c|c|c|c|c|c|}{
 \hline
 Estándar 802.11&Frecuencia&Velocidad máxima&Throughput típico&Alcance interior&Alcance exterior\\
 \hline
 a& 5 GHz& 54 Mbps& 23 Mbps& 35 m& 120 m\\
 b& 2.4 GHz& 11 Mbps& 4.3 Mbps& 35 m& 140 m\\
 g& 2.4 GHz& 54 Mbps& 19 Mbps& 38 m& 140 m\\
 n& 2.4 GHz& 150 Mbps& 74 Mbps& 70 m& 250 m\\
 n& 5 GHz& 150 Mbps& 74 Mbps& 70 m& 250 m\\
 \hline
}
 


%TODO http://www.extremetech.com/electronics/220355-next-generation-wi-fi-802-11ah-announced-with-almost-double-the-range-lower-power

%TODO http://www.speedguide.net/faq/what-is-the-actual-real-life-speed-of-wireless-374


\subsection{Modos de funcionamiento de las interfaces WLAN}
El estándar 802.11 define siete modos de operación de las interfaces WLAN. De éstos, los más corrientes son \textbf{modo Master}, \textbf{modo infraestructura o managed mode} y modo \textbf{Ad Hoc}. Las interfaces que se encuentran en modo Infraestructura únicamente pueden comunicarse con otra en modo Master; las que se encuentran en modo Ad Hoc, únicamente con otras interfaces en el mismo modo.

\begin{description}
    \item[Modo Master]
    La interfaz en modo Master funciona como punto de acceso (\textbf{Access Point, o AP}), y todo el tráfico entre las demás estaciones pasa a través de ella. 
    \item[Modo Infraestructura o Managed]
    Las interfaces en modo Managed se apoyan en la presencia de una interfaz en modo master (el AP), ante la cual deben registrarse (\textbf{asociarse}) para que el AP las tenga en cuenta.
    \item[Modo Independiente o Ad Hoc]
    Las interfaces se comunican directamente unas con otras sin necesidad de un punto de acceso, en forma \textbf{peer-to-peer}.
    \item[Modo Trama o Mesh]
    En topologías donde no todas las estaciones se encuentran en la zona de cobertura de los AP, las interfaces en modo Mesh funcionan como sistemas intermedios, que rutean en forma inteligente, y reenvían, los frames de las estaciones lejanas hacia los AP. 
    \item[Modo Promiscuo]
    La interfaz es capaz de leer todo el tráfico, independientemente de que la dirección destino de los frames sea o no igual a su dirección MAC. 
    \item[Modo Monitor]
    La interfaz es capaz de leer todo el tráfico, sin necesidad de estar asociada a un AP ni a una red Ad Hoc.
    \item[Modo Repetidor o Repeater]
    Una interfaz en modo Repetidor funciona como un bridge entre los APs de una red de infraestructura. Extiende el área de cobertura de la red y permite el \textit{roaming} o movilidad entre APs de una misma comunidad.
\end{description}



\subsection{Terminología IEEE 802.11}
Los elementos que componen las redes WiFi de tipo infraestructura (Fig. \ref{fig:infra}) tienen unos ciertos nombres precisos según los estándares de la familia 802.11.

\figura{infra} {Red de tipo Infraestructura} {infrastructure-crop.pdf}

\begin{description}
    \item[Medio inalámbrico] Es el medio utilizado para la transmisión (en las redes WLAN, el espacio). 
    \item[Estación (Station o STA)] Las estaciones son los hosts o dispositivos provistos de interfaces inalámbricas que forman parte de la red.
    \item[Punto de acceso (Access Point o AP)] Es una estación cuya interfaz funciona en un modo especial, bajo el cual no transmite datos propios de la estación, sino que recibe y redistribuye los frames de las demás estaciones. El papel del AP en las redes WiFi es análogo al de un \textbf{hub} Ethernet.    
    \item[Sistema de Distribución (Distribution System o DS)] Es la infraestructura de red cableada a la cual se pretende conectar la red WiFi. El AP posee una interfaz cableada, conectada al DS, y actúa como \textbf{bridge} entre la red WiFi y el DS.
\end{description}


\subsection{Tipos de tramas}
Los diferentes tipos de \textbf{tramas} o \textbf{frames} 802.11 cumplen diferentes objetivos dentro del complejo funcionamiento de las redes WLAN. 
\subsubsection{Tramas de administración}
Estos frames se utilizan para iniciar o terminar conexiones a una red de infraestructura. 
    \begin{itemize}
        \item Los frames de \textbf{beacon} o \textbf{baliza} son emitidos periódicamente por el AP. Permiten a las estaciones saber qué redes se encuentran disponibles cerca, y con qué características, para poder elegir a cuál red se unirán.
        \item Los frames de \textbf{Probe Request y Probe Response} tienen el mismo propósito, pero el descubrimiento de la red es iniciado por la estación.
        \item Los frames de autenticación sirven a las estaciones que desean asociarse al AP como primer paso del procedimiento.
        \item Los frames de asociación y des-asociación son emitidos por las estaciones para que el AP administre su pertenencia a la red.
    \end{itemize}  

\subsubsection{Tramas de control}
    \begin{itemize}
        \item Una estación confirma los frames recibidos con un frame de control especial, de reconocimiento, \textbf{ACK o Acknowledgement}. 
        \item Los frames cuya longitud se encuentra por encima de un cierto valor corren mayor riesgo de colisiones, por lo cual para esos casos se utiliza el procedimiento de Evitación de Colisiones (\textbf{Collision Avoidance o CA}) basado en una solicitud de permiso para transmitir (frame de \textbf{Request To Send o RTS} y un otorgamiento de dicho permiso (frame de \textbf{Clear To Send o CTS}). RTS y CTS son tramas de control.
    \end{itemize}  
\subsubsection{Tramas de datos}
    \begin{itemize}  
        \item Las tramas de datos tienen un formato inspirado en el de los frames Ethernet, que llevan direcciones MAC origen y destino. Sin embargo, en las redes de tipo infraestructura una estación utiliza un AP para llegar a otra estación, la cual puede formar parte del área de cobertura inalámbrica o estar directamente conectada al sistema de distribución cableado, y debido a esta particular forma de operación, las tramas 802.11 llevan tres (y a veces cuatro) direcciones MAC. 
    \end{itemize}

\subsection{Operación de 802.11}

\subsubsection{Ancho de banda compartido}
Analogía con hubs Ethernet
Atenuación
\subsubsection{Direccionamiento}
Mismas MAC que Eth
\subsubsection{Codificación}
Velocidades de transmisión variables
\subsubsection{Canales}
Interferencia y separación
\figura{canales} {Canales estandarizados por 802.11 en la frecuencia de 2.4 GHz} {canales-crop.pdf}
\subsubsection{Extensión del BSS}
WDS

\subsubsection{Tuning}


\begin{table}
\scriptsize
\parindent=0pt
\hrulefill
\begin{tabular}{   |p{\dimexpr 0.16\linewidth-2\tabcolsep} |
                   p{\dimexpr 0.28\linewidth-2\tabcolsep} |
                   p{\dimexpr 0.28\linewidth-2\tabcolsep} |
                   p{\dimexpr 0.28\linewidth-2\tabcolsep} | }
\hline
Parámetro 	& Significado  & Cuando se disminuye & Cuando se aumenta \\
\hline
\texttt{Beacon Interval}&Tiempo entre transmisión de beacons&Los scans pasivos se completan antes, y las estaciones se mueven más rápido sin perder conectividad
& Aumentan la capacidad de radio, el throughput y la vida de las baterías\\
\hline
\texttt{Fragmentation Threshold}&Los frames más largos que el FT se transmiten usando el procedimiento de fragmentación&En zonas con ruido, la interferencia solamente corrompe los fragmentos, aumentando el throughput&En zonas sin ruido, aumenta el throughput al reducirse la fragmentación \\
\hline
\texttt{RTS Threshold}&A los frames más largos que el RTS Threshold se les antepone el intercambio RTS/CTS&Mejor throughput donde hay muchas situaciones de nodos ocultos&Mejor throughput si no hay interferencia\\
\hline
\texttt{Listen Interval}&Las estaciones que entran en ahorro de energía despiertan al transcurrir esta cantidad de intervalos entre beacons del AP&Los frames unicast hacia la estación tienen menor latencia y la carga en memoria de los AP se ve reducida&El transceiver queda apagado más tiempo y así se conservan más tiempo las baterías de las estaciones móviles\\
\hline
\texttt{Long Retry Limit}&Cantidad de intentos de retransmisión para frames más largos que el RTS Threshold&Los frames se descartan más rápidamente y el espacio de memoria requerido es menor&La retransmisión lleva más tiempo y puede causar que el TCP disminuya la tasa de transferencia \\
\hline
\texttt{Short Retry Limit}&Cantidad de intentos de retransmisión para frames más cortos que el RTS Threshold&Lo mismo que para Long Retry Limit&Lo mismo que para Long Retry Limit \\
\hline
\end{tabular}
\caption{Algunos parámetros configurables en interfaces 802.11}
\label{tab:tuning}
\end{table}


\section {Estudio de caso I}


Una organización desarrolla sus actividades en un campus con dos edificios principales, que alojan las oficinas y talleres. En la organización existen tres áreas principales: Operaciones, Comercialización, e Ingeniería. 

Cada área utiliza un servidor de base de datos principal que es propio del área. La organización cuenta además con un servidor de web y de correo electrónico, un servidor de archivos y un servidor de backups, los tres de uso general para las tres áreas. Se desea que todos los puestos de trabajo puedan acceder además a Internet. 

Una preocupación especial de la organización es darle protección a los servidores y puestos de trabajo de Ingeniería, que no deben ser accedidos desde las demás áreas.
 
La planta del campus y sus edificios, con los principales puntos donde se necesita conectividad, es como indica la Fig. \ref{fig:caso01}. En este diagrama se muestran, con diferentes colores, los puestos de trabajo de cada área. 


\figura[12]{caso01}{Conectando el campus de una organización}{caso01.jpg}


En base a esta información, indique:
\begin{itemize}
	\item Dónde situaría los servidores mencionados.
	\item Dónde ubicaría los elementos de conectividad (switches y routers).
	\item Qué cantidad de puertos debería tener cada uno de estos elementos.
	\item Con qué medios (cobre, fibra, inalámbricos) vincularía los clientes y los elementos de conectividad.
	\item Cómo distribuiría direcciones IP para los clientes y servidores. 
	\item Si utilizaría alguna arquitectura de VLANs, cuál, y por qué. En caso afirmativo, cómo relacionaría las VLANs entre sí, qué enlaces entre switches y routers deben ser de trunking y por qué.
	\item Si utilizaría alguna forma de redundancia, y en caso afirmativo, cuál debe ser el camino normal del tráfico y por qué.
\end{itemize}








\newpage
\part {Redes Privadas Virtuales}

\section{Elementos de las VPN}
\begin{itemize}
	\item VPN = \emph{Virtual Private Network}
	\item Un software específico (cliente-servidor) establece una red \emph{virtual} entre las entidades
	\item Sólo aquellos miembros autorizados de la red virtual \emph{privada} pueden utilizarla
	\item \emph{Tunneling} o encapsulamiento de ciertas capas en otras
	\item Encriptación
\end{itemize}


\subsection{Motivación de las VPN}
\begin{itemize}
	\item El software de la organización evoluciona a aplicaciones colaborativas (Groupware, CRM , ERP...)
	\item Trabajo de todos los colaboradores sobre el espacio digital de la organización
	\item Diferentes sedes, o colaboradores (domiciliarios o móviles), en diferentes lugares de la región o del mundo
	\item Soporte universal de Internet en lugar de enlaces arrendados, dedicados
	\item Requerimientos de seguridad
	\item Firewalls y NAT requieren \emph{conduits} o perforaciones para acceder a los servicios de las sedes
 \end{itemize}



\subsection{Modelos de VPN}
\begin{itemize}
	\item Implementadas y administradas por un proveedor de comunicaciones
	\item Implementadas y administradas por el usuario sobre un transporte común (Internet)
\end{itemize}


\subsection{Niveles de implementación}
\begin{itemize}	 

	\item De nivel de Enlace o de capa 2

	\begin{itemize}
		\item Único espacio IP y dominio de broadcast
		\item La VPN trafica frames
		\item Bridging
	\end{itemize} 

	\item De nivel de Red o de capa 3

	\begin{itemize}
		\item  Varias redes con diferentes espacios IP y un servidor VPN que las conecta como backbone 
		\item La VPN trafica paquetes
		\item Ruteo
	\end{itemize}
\end{itemize}


\subsection{Algunas tecnologías}
\begin{itemize}
	\item GRE (RFCs 1701, 1702) Mecanismo general de tunneling, no define encriptación
	\item Protocolos de capa 2: PPTP (Microsoft), L2F, L2TP, L2Sec
	\item Protocolos de capa 3: IPSec
	\item Protocolos de capa 4: SSL/TLS, OpenVPN
\end{itemize}


\figura{vpnbrid} {VPN de capa 2, equivalente a un bridge} {vpn1-bridge.jpg}

\figura{vpnrut} {VPN de capa 3, equivalente a un router} {vpn1-router.jpg}


\section{OpenVPN}

\begin{itemize}
	\item OpenVPN puede funcionar como VPN de capa 2 o de capa 3. 
	\item Proyecto libre compatible con Windows.
	\item No interoperable con IPSec.
	\item Dos procesos, cliente y servidor, corren en diferentes hosts en diferentes redes. 
	\item Cliente y servidor establecen una conexión TCP/IP que utilizan como túnel.
	\item El resto del trabajo lo hace el bridging (en las VPN de capa 2) o el ruteo (de capa 3). 
	\item La arquitectura de OpenVPN (y otras implementaciones de redes virtuales) es un ejemplo claro de uso de interfaces virtuales como los bridges y tun/tap. 
\end{itemize}

\subsection{Soporte necesario para OpenVPN}
\begin{description}
	\item [Tun/Tap] Dispositivos de red virtuales. Trafican unidades de datos entre procesos. 
	\begin{itemize}
		\item Un dispositivo tun (de capa 3) trafica paquetes, mientras que un tap (de capa 2) intercambia frames. 
		\item El propósito básico de Tun/Tap es la creación de túneles.
		\item Un proceso que tiene un dispositivo tun abierto puede escribir un paquete en él.  Del mismo modo, un dispositivo tun que recibe un paquete IP lo envía a un proceso de usuario que tiene abierto ese dispositivo.
		\item Un proceso que tiene un dispositivo tap abierto puede escribir un frame en él.  Del mismo modo, un dispositivo tap que recibe un frame lo envía a un proceso de usuario que tiene abierto ese dispositivo.
	\end{itemize}
	\item [Bridge] Un dispositivo bridge de software funciona exactamente igual que un bridge físico pero entre interfaces del mismo host. Es decir, copia frames de una a otra interfaz de capa 2 (reales o virtuales) efectuando filtrado por dirección MAC e inundando los broadcasts.
	\begin{itemize}
		\item Un bridge Linux \emph{esclaviza} a dos o más interfaces (reales o virtuales).
		\item Al ser incorporada una interfaz a un bridge, pierde su dirección IP y ésta pasa al bridge. 
		\item Los bridges Linux se controlan desde la línea de comandos con el comando brctl.
\end{itemize}
	\item[Ruteo] El kernel utiliza la tabla de ruteo para mover paquetes entre interfaces. Cuando aparece por una interfaz (real o virtual) un paquete de capa 3, se consulta la tabla de ruteo, que dice en cuál interfaz (real o virtual) debe copiarse.  
	\item [OpenSSL] Biblioteca de encriptación de uso general, para proteger la privacidad del tráfico a través de la VPN.
	\item [LZO] Biblioteca de compresión de datos para mejorar la performance del túnel creado por la VPN.
\end{description}



\subsection{Funcionamiento de OpenVPN en capa 2}

\begin{enumerate}
	\item En el host cliente:
	\begin{itemize}
		\item Se crea un dispositivo tapX.
		\item Se establece un bridge entre una interfaz local (ethX) y el tapX.
		\item La dirección IP de la interfaz local pasa al bridge.
		\item El proceso cliente de OpenVPN abre el dispositivo tapX y además establece conexión con el servidor a través de Internet.
		\item Cuando el host cliente reciba de su red local un frame por ethX, el bridge filtrará el frame hacia el tapX. El proceso cliente de OpenVPN recibirá el frame y lo derivará por la conexión TCP/IP hacia el servidor.
		\item Cuando el proceso cliente reciba un frame a través de la conexión TCP/IP, lo escribirá por la interfaz tapX. El bridge del host cliente recibirá el frame y le aplicará la acción de bridging correspondiente. 
	\end{itemize}
	\item En el host servidor:
	\begin{itemize}
		\item Se crea un dispositivo tapX.
		\item Se establece un bridge entre una interfaz local (ethX) y el tapX.
		\item La dirección IP de la interfaz local pasa al bridge.
		\item El proceso servidor de OpenVPN abre el dispositivo tapX y espera pedido de conexión del cliente a través de Internet.
		\item Cuando el host servidor reciba de su red local un frame por ethX, el bridge filtrará el frame hacia el tapX. El proceso servidor de OpenVPN lo tomará y lo derivará por la conexión TCP/IP hacia el cliente.
		\item Cuando el proceso servidor reciba un frame a través de la conexión TCP/IP, lo escribirá por la interfaz tapX. El bridge del host cliente recibirá el paquete y le aplicará la acción de bridging correspondiente. 
	\end{itemize}
\end{enumerate}

\subsection{Funcionamiento de OpenVPN en capa 3}

\begin{enumerate}
	\item En el host cliente:
	\begin{itemize}
		\item Se crea un dispositivo tunX.
		\item El proceso cliente de OpenVPN abre el dispositivo tunX y además establece conexión con el servidor a través de Internet.
		\item El tunX local recibe dirección IP A. Queda establecido un túnel entre procesos cliente y servidor cuyos extremos tienen direcciones IP A y B.
		\item Se establece una ruta que dirige los paquetes destinados a la red Y via tunX (próximo salto: B).
		\item Cuando el host cliente reciba de su red local, o de procesos locales, un paquete destinado a la red Y, lo escribirá en la interfaz de red virtual tunX. El proceso cliente de OpenVPN lo recibirá y lo derivará por la conexión TCP/IP hacia el servidor.
		\item Cuando el proceso cliente reciba un paquete a través de la conexión TCP/IP, lo escribirá por la interfaz tunX. El kernel del host cliente recibirá el paquete y le aplicará las reglas de ruteo correspondientes. 
	\end{itemize}
	\item En el host servidor:
	\begin{itemize}
		\item Se crea un dispositivo tunX.
		\item El proceso servidor de OpenVPN abre el dispositivo tunX y espera pedido de conexión del cliente a través de Internet.
		\item El tunX local recibe dirección IP B. Queda establecido un túnel entre procesos cliente y servidor cuyos extremos tienen direcciones IP A y B. 
		\item Se establece una ruta que dirige los paquetes destinados a la red X via tunX (próximo salto: A).
		\item Cuando el kernel del host servidor reciba de su red local, o de procesos locales, un paquete destinado a la red X, lo escribirá en la interfaz de red virtual tunX. El proceso servidor de OpenVPN tomará el paquete y lo derivará por la conexión TCP/IP hacia el cliente.
		\item Cuando el proceso servidor reciba un paquete a través de la conexión TCP/IP, lo escribirá por la interfaz tunX. El kernel del host servidor recibirá el paquete y le aplicará las reglas de ruteo correspondientes. 
	\end{itemize}
\end{enumerate}

\figura{vpnl2} {VPN de capa 2, trafica frames} {vpn-l2.jpg}

\figura{vpnl3} {VPN de capa 3, trafica paquetes} {vpn-l3.jpg}

\section{Configuración de OpenVPN}

\subsection{Autenticación mediante clave privada}
\begin{itemize}
	\item PKI = Public Key Infrastructure
	\item La seguridad en redes comprende problemas como la confidencialidad y la integridad de los mensajes, y la autenticación de los usuarios o dispositivos
	\item Clave simétrica, con secreto compartido, vs. Claves asimétricas
	\item Dificultad de la distribución de claves, solución: claves asimétricas
	\item Cada usuario tiene $K+$ = Clave pública, $K-$ = Clave privada
	\item Dos escenarios de uso:
	\begin{itemize}
	\item Cuando A encripta con la $K+$ de B:
	\begin{itemize}
		\item A envía a B en forma segura, nadie puede leer el mensaje
		\item Sólo B puede desencriptarlo con su $K-$
		\item Se asegura la confidencialidad
	\end{itemize}
	\item Cuando A encripta con su propia $K-$:
	\begin{itemize}
		\item Cualquiera puede leer el mensaje con la $K+$ de A
		\item Pero sólo A puede haberlo escrito, con su $K-$ (firma digital)
		\item Se aseguran la integridad y la autenticidad
	\end{itemize} 
	\end{itemize}
	\item El problema al utilizar la $K+$ de A que hemos recibido es \emph{asegurar que es realmente la de A}
	\item Se resuelve mediante autoridades de certificación (CA), entidades confiables
	\item La clave pública de A, junto con la identidad de A, firmada digitalmente por una CA, es un \emph{certificado} de A
	\item Conexión segurizada por PKI
	\begin{enumerate}
		\item Fase de autenticación mutua mediante intercambio de certificados, validando los certificados según la CA.
		\item Usando las $K+$ recibidas, las partes acuerdan en forma segura una clave simétrica que sirve para esta sesión.
		\item El resto de la comunicación se encripta usando esta clave simétrica (métodos DES, 3DES, AES).
	\end{enumerate} 
\end{itemize}


\subsection{Crear una Autoridad de Certificación}

En el servidor de VPN, copiar el directorio \lstinline{/usr/share/doc/openvpn/examples/easy-rsa} y todos sus contenidos a /etc/openvpn
\begin{lstlisting}
cd /etc/openvpn/easy-rsa/2.0
vi vars
. vars
./clean-all
./build-dh
./build-ca
./build-key-server servidor
./build-key cliente1
./build-key cliente2
# Migrar los certificados
# En el server:
# ca.crt ca.key dh1024.pem r3.crt r3.key
# 
# En el cliente:
# ca.crt r1.crt r1.key
\end{lstlisting}


\subsection{Documentación online}
\begin{itemize}
	\item Proyecto Lihuén\footnote{\url{http://lihuen.linti.unlp.edu.ar/index.php?title=Configurando_Redes_Privadas_Virtuales_con_OpenVPN}}
	\item OpenVPN HOWTO\footnote{\url{http://openvpn.net/index.php/open-source/documentation/howto.html}}
\end{itemize}


\subsection{Vinculación de LANs mediante Openvpn}
\begin{itemize}
	\item Instalar OpenVPN en los nodos de frontera de las LANs.
	\item Decidir cuál nodo será el servidor. Debe tener dirección IP pública accesible desde los clientes. No es necesaria la traducción DNS salvo que se trate de una IP dinámica.
	\item Preparar una autoridad de certificación (por ejemplo, en el servidor), y certificados para el servidor y para los clientes.
	\item Se pueden utilizar los modelos de archivos de configuración existentes en \url{/usr/share/doc/openvpn/examples}.
	\item En el directorio \lstinline{/etc/openvpn} del servidor preparar el archivo de configuración \lstinline{server.conf}, y en los clientes, \lstinline{client.conf}. Al arrancar, OpenVPN lee todos los archivos con extensión .conf que existan en ese directorio, y crea un proceso por cada uno, con la configuración que contengan.
	\item Modificar los parámetros de configuración:
	\begin{enumerate}
		\item Cliente o servidor
		\item Dirección y port, nodo local o remoto
		\item Protocolo UDP o TCP
		\item Device Tun o Tap, según la capa de operación de la VPN  
		\item Nombres de los archivos de clave y certificado 
		\item Archivo de log, útil para debugging de la configuración y operación
		\item Rutas a redes propias que deban ser inyectadas en el peer
\end{enumerate} 
	\item Arrancar o detener la VPN con el comando \lstinline{/etc/init.d/openvpn [start|stop]}.
\end{itemize}

\subsubsection {OpenVPN de nivel de Red}
Al configurar OpenVPN en capa 3, o nivel de Red, la configuración de OpenVPN puede asumir la creación de rutas entre las redes clientes y las de detrás del servidor. Estas rutas aparecen y desaparecen en la tabla de ruteo de los nodos extremos de la VPN según se activa o desactiva el proceso de OpenVPN.

\begin{itemize}
	\item Para que el cliente conozca las redes detrás del servidor, éste le inyecta las rutas correspondientes con la opción push de la configuración.
	\item El servidor instala rutas a las redes detrás de los clientes si se  especifican con la directiva route de la configuración del servidor. Además, para cada cliente, debe haber un archivo con la directiva iroute dentro del subdirectorio /etc/openvpn/ccd.
\end{itemize}

\subsubsection {OpenVPN de nivel de Enlace}
Al configurar OpenVPN en capa 2, o nivel de Enlace, la configuración debe incluir la administración del bridge entre la interfaz de red local y el dispositivo Tap que es utilizado por la VPN. Los frames que llegan a la interfaz de la red local serán copiados por el bridge en el Tap y seguirán su camino a través de la VPN.

Para la administración del bridge hay dos estrategias básicas: 
\begin{itemize}
	\item El bridge puede tener una definición estática, y ser siempre activado cada vez que arranca el equipo, y sólo ser desactivado cuando se detiene el equipo.
	\item El bridge puede ser activado y desactivado cuando se inicia y detiene OpenVPN. 
\end{itemize}


En el primer caso, la existencia del bridge es completamente independiente de la activación de la VPN. Para esta opción existe una forma de configuración del bridge que es dependiente del sistema operativo (ver Anexo \ref{subsec:staticbridge}). 

Para la segunda opción, es conveniente utilizar los scripts de creación de tap y montado de bridges que se muestran en el Anexo \ref{sec:bridgeupdown} y dispararlos con las directivas up y down de la configuración de OpenVPN. Por ejemplo:

\begin{lstlisting}
...
ca ca.crt
cert server.crt
key server.key
dh dh1024.pem
...
up /etc/openvpn/bridge-up
down /etc/openvpn/bridge-down
\end{lstlisting}


\section{Temas de práctica}
El laboratorio de la Fig. \ref{fig:vpn-lab1} se encuentra implementado sobre máquinas virtuales. 
\begin{comment}Utilizando una consola, encontrará en el directorio \lstinline{labs/openvpn-1} una configuración de laboratorio que se arranca con el comando \lstinline{lstart} y se detiene con \lstinline{lhalt}. 
\end{comment}


\figura{vpn-lab1}{Laboratorio 1 OpenVPN}{vpn-lab1.jpg}

El equipo virtual $R2$, que es gateway default de los otros routers, simula la Internet, en el sentido de que descarta todo tráfico dirigido a redes privadas. Por este motivo las redes con prefijo 10.0.0.0/8, del laboratorio, no son accesibles una desde la otra, siendo necesaria una solución de Red Privada Virtual. La configuración de este equipo $R2$ no debe ser modificada.

\begin{itemize}
	\item Verifique que las direcciones, ruteo default y topología corresponden a la figura.
	\item Verifique servicio de nombres, ping y traceroute a www.google.com.
	\item Ping a los routers de la topología.
	\item Ping a los hosts de la misma LAN.
	\item Ping a los hosts de la LAN opuesta. 
	\item Compruebe navegación en ambiente de caracteres con el comando \texttt{lynx http://www.google.com}.
	\item Compruebe que puede correr el web server Apache en cualquiera de las PCs. 
	\item Compruebe navegación en ambiente de caracteres desde una PC al servidor Apache del otro nodo de la misma LAN.
\end{itemize}

Una vez que esté familiarizado con el escenario:
\begin{enumerate}
	\item Implementar una VPN de capa 3 entre ambas LANs. Verifique las interfaces existentes en los nodos cabecera y el contenido de sus tablas de ruteo. Comprobar que los  clientes de una de las redes pueden utilizar recursos (ssh, servidor http) de clientes de la otra. Instale servicios basados en broadcast (como Samba, sobre el protocolo SMB) y observe desde qué lugares pueden accederse.
	\item Modificar el direccionamiento IP del laboratorio para que todos los clientes reciban direcciones en la misma subred e implementar una VPN de capa 2. Comprobar que ambas redes forman un único dominio de broadcast. Verifique el comportamiento de los protocolos basados en broadcast como SMB. 
	\item ¿Usaría transporte UDP o TCP para una VPN en capa 2?
\end{enumerate}


\newpage
\part {Balance de carga  y Alta Disponibilidad  en redes}
\section{Bonding}


Acoplamiento de dos o más interfaces de red, creando una interfaz virtual capaz de funcionar en diferentes modos. Se conoce con diferentes nombres (\textit{channel bonding, teaming, link aggregation}). 
Los modos de configuración de bonding definen el comportamiento del conjunto de interfaces y cumplen con diferentes objetivos. Algunos modos proporcionan tolerancia a fallos mediante redundancia; otros aumentan el ancho de banda disponible por agregación de enlaces.

Los modos active-backup, balance-tlb, y balance-alb no requieren una configuración especial en los switches a los cuales está conectado el bond.
Sin embargo, los modos 802.3ad, balance-rr, balance-xor y broadcast requieren capacidades especiales del switch definidas en documentos IEEE.


\begin{description}
\item [Modo 0 (balance-rr)]
Este modo transmite frames en orden secuencial desde el primer esclavo disponible hasta el último. Si un bond tiene dos interfaces reales, y llegan simultáneamente dos frames a ser enviados desde la interfaz bond, el primero será transmitido por el primer esclavo; el segundo frame, por el segundo esclavo; el tercer frame será transmitido por el primer esclavo, etc. Esto provee a la vez balance de carga y tolerancia a fallos. 

\item [Modo 1 (active-backup)]
Este modo coloca a las interfaces esclavas en estado de backup, y sólo se activará una de ellas si se pierde el link de la interfaz activa. En este modo, sólo hay un esclavo activo en el bond en cada momento. Sólo se activa un esclavo diferente si falla el esclavo activo. Este modo provee tolerancia a fallos.

\item [Modo 2 (balance-xor)]
Transmite basándose en una fórmula XOR. Se computa la operación XOR entre la dirección MAC de origen y la de destino, módulo la cantidad de esclavos (es decir, se toma el resto de dividir por la cantidad de esclavos). Este procedimiento tiene el efecto de seleccionar siempre el mismo esclavo para cada dirección MAC destino. Provee balance de carga y tolerancia a fallos.

\item [Modo 3 (broadcast)]
Este modo transmite todos los frames por todas las interfaces esclavas. Es el menos usado, sólo para propósitos específicos, y sólo provee tolerancia a fallos. 

\item [Modo 4 (802.3ad)]
Este modo se conoce como el modo de agregación dinámica de enlaces (Dynamic Link Aggregation). Crea grupos de agregación que comparten la misma velocidad y modos de duplexing. Este modo requiere un switch que soporte la norma IEEE 802.3ad (Dynamic Link).

\item [Modo 5 (balance-tlb)]
Llamado balance de carga adaptativo en transmisión. El tráfico de salida se distribuye de acuerdo a la carga actual y es encolado en cada interfaz esclava. El tráfico entrante es recibido por el esclavo actual. 

\item [Modo 6 (balance-alb)]
Este modo es el de balance de carga adaptativo. Esto incluye balance-tlb y balance de carga en recepción (rlb) para tráfico IPv4. El balance de carga en recepción se logra por negociación ARP. El driver de bonding intercepta las respuestas ARP enviadas por el servidor y sobreescribe la dirección MAC origen con la dirección MAC única de uno de los esclavos en el bond, de forma que diferentes clientes usen diferentes direcciones MAC para dirigirse al server. 
\end{description}

\subsection {Detección de eventos}
Los modos que ofrecen tolerancia a fallos necesitan algún mecanismo para detectar eventos de caída de las interfaces de red (NIC) locales, los enlaces, o las NICs de los extremos opuestos de los vínculos. Ante la detección de un evento de fallo, el bond fuerza la conmutación a otra interfaz, que entonces se convierte en primaria. Esta acción se llama \emph{failover}.  

Para la detección de eventos hay dos opciones posibles.

\begin{enumerate}
		\item MII (Medium Independent Interface). Especificación que cumplen la mayoría de las NICs modernas, que presenta datos de link activo o inactivo en forma independiente de la implementación del medio conectado. Se debe especificar los parámetros \lstinline$bond-miimon$ (intervalo de revisión del estado del link), \lstinline$bond-downdelay$ (tiempo desde que se detecta fallo hasta que se da de baja la interfaz) y \lstinline$bond-updelay$ (tiempo para volver a poner en servicio la interfaz una vez que vuelve el link a estado activo). 
		\item ARP. Se establece por configuración una cantidad de direcciones IP de control en la red local, y el bond emite consultas ARP periódicas a estas direcciones. Cuando se deja de recibir respuesta por la interfaz activa durante una cantidad de tiempo, configurable, se considera que ha caído el enlace y se realiza el failover. Se deben configurar los parámetros \lstinline$bond-arp-interval$ (intervalo entre emisión de ARP) y \lstinline$bond-arp-ip-target$ (lista de IPs confiables).
\end{enumerate}
  
\subsection {Configuración en Debian}

\subsubsection {Configuración con MII}
\begin{lstlisting}
# /etc/network/interfaces
# The loopback network interface
auto lo
iface lo inet loopback
auto eth0
iface eth0 inet manual
	bond-master bond0
	bond-primary eth0 eth1

auto eth1
iface eth1 inet manual
	bond-master bond0
	bond-primary eth0 eth1

auto bond0
iface bond0 inet static
	address 192.168.1.15
	netmask 255.255.255.0
	network 192.168.1.0
	gateway 192.168.1.1
	bond-slaves eth0 eth1
	bond-mode active-backup
	bond-miimon 100
	bond-downdelay 200
	bond-updelay 200

\end{lstlisting}

\subsubsection {Configuración con ARP}

\begin{lstlisting}
# /etc/network/interfaces
# The loopback network interface
auto lo
iface lo inet loopback

# No se especifican las auto ethX

auto bond0
iface bond0 inet static
	address 10.1.1.1
	netmask 255.255.255.0
	network 10.1.1.0
	bond_primary eth0
	slaves eth0 eth1
	bond-mode active-backup
	bond-arp-interval 2000
	bond-arp-ip-target 10.1.1.2 10.1.1.3

\end{lstlisting}



\subsection{Temas de práctica}

La topología de la Fig. \ref{fig:bonding} comprende un nodo llamado switch, y tres nodos conectados a él. El switch tiene cinco interfaces, cada una conectada a un enlace diferente. Dos de los nodos tienen dos interfaces cada uno, y el tercero una sola. Cada interfaz de los nodos host1 a host3 está conectada a su propio enlace. Los enlaces se denominan, de izquierda a derecha en el diagrama, link1 a link5. 

\figura[10]{bonding}{Configuración del laboratorio de bonding}{bonding.png}

\recuadro{

¡Ver las Notas sobre el laboratorio en la sección más abajo!

}


\begin{enumerate}
	\item Configure en forma estática la red del nodo denominado switch implementando un bridge que esclavice a todas sus interfaces. Dé direcciones en el mismo dominio de broadcast a los demás nodos. Compruebe llegada por ping. Observe el tráfico que pasa por las interfaces del switch con el comando tcpdump -i ethX. El nodo virtual switch, ¿resulta un buen modelo de un switch de hardware? ¿Necesita tener una dirección IP? 
	\item En los nodos host1 y host2, Instale el módulo que permite esclavizar las interfaces con \code{dpkg -i /hostlab/ifenslave-...deb}.
	\item Configure el nodo host1 indicando en /etc/modules que debe ser cargado el módulo bonding al arranque. Configure la red del mismo nodo, ligando ambas interfaces mediante un bond en modo active-backup. ¿Qué capacidades tiene este modo?
	\item Mantenga un ping desde host3 a host1, observando el tráfico por las interfaces del switch. Simule la caída del vínculo con host1.  ¿Se interrumpe el ping? ¿Sigue pasando el tráfico por las mismas interfaces?
	\item Configure el nodo host2 del mismo modo que en el punto 3, pero con el bond en modo balance-rr. ¿Qué capacidades tiene este modo?
	\item Ejecute el mismo experimento del punto 4 pero desde host3 a host2. 
\end{enumerate}


\subsubsection{Notas sobre el laboratorio}
\begin{itemize}
	\item Es necesario instalar el paquete \code{ifenslave}. En nuestro laboratorio se ha provisto un ejemplar del archivo DEB correspondiente en el directorio \code{/hostlab}.
	\item Es necesario indicar el modo, técnica de detección de eventos, y parámetros, en el archivo \code{/etc/modules} (a pesar de indicarlo en \code{/etc/network/interfaces}), de la siguiente manera:
	\begin{lstlisting}
	alias bond0 bonding
	options bonding mode=1 miimon=100 downdelay=200 updelay=200
	\end{lstlisting}
(detección por MII) o bien:
	\begin{lstlisting}
	alias bond0 bonding
	options bonding mode=1 arp_interval=1000 arp_ip_target=10.1.1.1
	\end{lstlisting}
(detección por ARP).
	\item Al probar sistemas de Alta Disponibilidad, una parte muy importante es la inyección de fallas. La tecnología de virtualización utilizada para el laboratorio no permite replicar correctamente el evento de caída de interfaz, enlace o peer, cuando se usa la técnica de detección de fallas MII (la interfaz aparece siempre conectada). La mejor aproximación que hemos logrado a la simulación de la falla es cuando se elige detección de fallas por ARP.
	\item Para simular la caída de un vínculo se sugiere detener el proceso UML que simula el enlace. Logramos esto buscando entre todos los procesos \code{uml_switch} de la máquina host aquel relacionado con el nombre del enlace. Por ejemplo, para simular la caída del link1:
	\begin{lstlisting}
# ps f | grep uml_switch
23137 pts/1    S+     0:00  \_ grep uml_switch
...
19381 pts/1    S      0:00 /home/rii/netkit/bin/uml_switch -hub -unix /home/oso/.netkit/hubs/vhub_oso_link1.cnct
# kill -STOP 19381
\end{lstlisting}
	El proceso se hace continuar (se restablece el vínculo virtual) con \code{kill -CONT 19381}.
	\item Cada vez que hagamos un cambio de configuración será preferible, antes de hacer una nueva prueba, bajar y volver a levantar el laboratorio completo con \code{lhalt} y \code{lstart}. 
	\item La interfaz virtual bond puede monitorearse mirando el pseudo archivo correspondiente en el directorio /proc:
\begin{lstlisting}
# cat /proc/net/bonding/bond0 
Ethernet Channel Bonding Driver: v3.2.5 (March 21, 2008)

Bonding Mode: fault-tolerance (active-backup)
Primary Slave: eth0
Currently Active Slave: eth0
MII Status: up
MII Polling Interval (ms): 0
Up Delay (ms): 0
Down Delay (ms): 0
ARP Polling Interval (ms): 2000
ARP IP target/s (n.n.n.n form): 10.1.1.2, 10.1.1.3

Slave Interface: eth0
MII Status: up
Link Failure Count: 0
Permanent HW addr: a6:8a:d7:63:1e:49

Slave Interface: eth1
MII Status: up
Link Failure Count: 0
Permanent HW addr: e6:c9:c8:a3:eb:72
\end{lstlisting}
\end{itemize}



\subsection{Referencias}
\begin{itemize}
	\item \url{https://www.kernel.org/doc/Documentation/networking/bonding.txt}
	\item \url{http://www.linuxfoundation.org/collaborate/workgroups/networking/bonding}
	\item \url{http://www.cyberciti.biz/tips/debian-ubuntu-teaming-aggregating-multiple-network-connections.html}
\end{itemize}



\section{Balance de Carga}


Consideraremos el problema de una organización con dos o más accesos a Internet a través de diferentes proveedores. Los enlaces extra se han contratado con la idea de seguir obteniendo acceso en caso de que uno de los proveedores falle (Fig. \ref{fig:split}). Sin embargo, resulta costoso diseñar una solución en modo \textit{activo-standby} (con sólo uno de los enlaces activo y sin aprovechar los demás), por lo cual se debe considerar alguna forma de distribución de tráfico por todos los enlaces. El ruteo default permite indicar solamente un camino de salida, de manera que se necesitará alguna técnica adicional para distribuir el tráfico. En general, cuando se necesita establecer alguna decisión administrativa sobre cómo se utilizan las rutas disponibles, se dice que estamos frente a una situación de ruteo por políticas. Una política es simplemente un conjunto de decisiones. 


\figura[11]{split}{Escenario de acceso dividido}{split}

Por la naturaleza de los protocolos de transporte de Internet, en general no será posible utilizar el total del ancho de banda agregado de todos los enlaces para una sola conexión, ya que segmentos TCP o UDP pertenecientes a la misma interacción con un servidor llegarán a destino provenientes de direcciones IP diferentes (aquellas que resulten del ruteo realizado por los diferentes proveedores). Este patrón de tráfico no puede formar parte de la misma conexión TCP o conversación UDP. Es decir, si se tienen dos accesos de 1Mbps cada uno, no es posible utilizar ambos enlaces agregados para efectuar una única transferencia de datos a 2Mbps en una sola conexión. En cambio, sí es posible distribuir la totalidad de la carga entre todos los vínculos, destinando parte de las conexiones por uno u otro enlace. 


El escenario de múltiples accesos a Internet suele llamarse \textit{Split Access} (acceso dividido). Los dos problemas de Split Access son:

\begin{enumerate}
	\item Hacer que el tráfico ingresado al sistema desde Internet por una interfaz vuelva por la misma interfaz.
	\item Lograr distribución o balance de carga.
\end{enumerate}

Las principales herramientas de software para conseguir estos objetivos son el paquete \texttt{iproute2} y el comando \texttt{iptables}. Con el primero definiremos una estructura de ruteo especial para el caso de uso de Split Access y con el segundo manipularemos el tráfico para aplicar diferentes políticas de ruteo.

\subsection{Paquete iproute2}
Paquete de configuración de red integral, que reemplaza a utilitarios como \texttt{ifconfig}, \texttt{route}, \texttt{netstat}. Permite configurar interfaces, información de ARP, ruteo por políticas, túneles, etc. Consiste de dos comandos: \texttt{ip}, que controla estado de interfaces y configuración de IPv4 e IPv6, y \texttt{tc}, que administra aspectos de calidad de servicio (\textit{Quality of Service}, QoS). Ejemplos en Cuadro \ref{tab:iproute2}.

\subsection{Comando iptables}
Es la interfaz de usuario al subsistema Netfilter del kernel, que permite modificar la forma como los paquetes IP atraviesan el kernel en su tránsito entre una interfaz de entrada, procesos locales, e interfaz de salida. Según el caso de cada paquete, aplica una serie de reglas impuestas por el usuario y agrupadas en diferentes \textbf{cadenas} (\texttt{PREROUTING, INPUT, OUTPUT, POSTROUTING, FORWARD}). Las reglas se almacenan en \textbf{tablas} que se aplican secuencialmente dentro de cada cadena. Las cadenas se disponen en diferentes lugares del recorrido del paquete por el kernel, y el usuario puede crear las suyas propias. 

Cada tabla tiene una aplicación específica:

\begin{description}
	\item[Tabla \texttt{filter}] Aplica reglas de filtrado para implementar políticas de firewalling. Es la tabla por defecto.
	\item[Tabla \texttt{nat}] Ejecuta edición de direcciones. Por ejemplo, cuando se aplica NAT (\textit{Network Address Translation}), convierte la dirección de salida de un paquete a una dirección propia.
	\item[Tabla \texttt{mangle}] Permite efectuar cualquier modificación de cualquier zona del paquete, además de afectar datos que acompañan al paquete durante su viaje por el kernel. 
	\item[Tabla \texttt{raw}] Implementa excepciones al seguimiento de conexiones.
	\item[Tabla \texttt{security}] Implementa control de acceso obligatorio y permite separar las acciones que puede tomar el usuario de las que impone el sistema.
\end{description}

Netfilter diferencia entre paquetes \textbf{destinados al host} (la dirección IP destino del paquete es alguna de las propias), \textbf{originados en el host} (la dirección origen es una dirección propia), o \textbf{en tránsito} (ni la dirección origen ni la dirección destino son propias). Este último caso es el que se presenta habitualmente en un router. 

\begin{description}
	\item[Ingreso] A todos los paquetes ingresantes, antes de que el proceso de ruteo determine si son o no dirigidos al host, se les aplican las reglas contenidas en las tablas \texttt{mangle} y \texttt{nat} de la cadena \texttt{PREROUTING}.
	\item[Egreso] Todos los paquetes que salen del host por cualquier motivo atraviesan las tablas \texttt{mangle} y \texttt{nat} de la cadena \texttt{POSTROUTING}. 
	\item[Paquetes destinados al host] Luego de determinarse que el paquete está destinado al host, éste atraviesa la cadena de \texttt{INPUT}, con las reglas que haya en la tabla \texttt{mangle} y la tabla \texttt{filter} de esa cadena, en ese orden (Fig. \ref{fig:iptables-input}). 
	\item[Paquetes originados en el host] Son producidos por procesos locales que emiten pedidos de servicio o responden a servicios solicitados. Atraviesan la cadena de \texttt{OUTPUT} (\texttt{mangle}, \texttt{nat} y \texttt{filter}, en ese orden) y luego se dirigen a la cadena de \texttt{POSTROUTING} (Fig. \ref{fig:iptables-output}). 
	\item[Paquetes que atraviesan el host] Se trata de paquetes que el router debe reenviar. Atraviesan las tablas \texttt{mangle} y \texttt{filter} de la cadena \texttt{FORWARD} y luego se dirigen a la cadena \texttt{POSTROUTING} (Fig. \ref{fig:iptables-forward}).
\end{description}


\tabla{iproute2}{Ejemplos de uso del comando ip}{l|l}{
\texttt{ip link list}					& Consultar interfaces \\
\texttt{ip address show}				& Consultar direcciones de interfaces \\
\texttt{ip neigh show}				& Tabla ARP\\
\texttt{ip route show}				& Rutas\\
\texttt{ip route list table main}		& Rutas, tabla principal\\
\texttt{ip route list table T}		& Rutas, tabla T\\
\texttt{ip route flush table main}	& Borrar rutas de tabla principal\\
\texttt{ip route flush cache}			& Borrar cache de rutas\\
\texttt{ip rule add from 10.0.0.0/24 table T}	& Agregar regla de asignación de tablas\\
\texttt{ip rule list}					& Consultar reglas de asignación\\
}
 
 
\figura[11]{iptables-input}{Recorrido de un paquete destinado al host}{iptables-input.pdf}
\figura[11]{iptables-output}{Recorrido de un paquete originado en el host}{iptables-output.pdf}
\figura[11]{iptables-forward}{Recorrido de un paquete reenviado por el host}{iptables-forward.pdf}

\subsection{Ruteo por políticas}
Un host determina por cuál de sus interfaces debe emitir un paquete mediante el proceso de ruteo, usando la información de rutas contenida en una tabla de ruteo o reenvío. Habitualmente, la tabla utilizada para el ruteo regular es única (la tabla \texttt{main}), pero con la técnica de ruteo por políticas, se pueden utilizar diferentes tablas de ruteo. 

Una vez seleccionada una tabla de ruteo, el kernel utilizará las rutas en dicha tabla para determinar qué interfaz de salida le corresponde a un paquete. Para cada tabla de ruteo se aplica siempre el mecanismo clásico de examinar la dirección destino del paquete y localizar en esa tabla el prefijo coincidente más largo. 

Ante cada paquete nuevo que llega, el proceso de ruteo buscará primero reglas para seleccionar una tabla de ruteo, y sólo si no las encuentra, aplicará las rutas en la tabla \texttt{main}. En particular, una ruta default en una tabla específica tendrá prioridad sobre la ruta default habitual. 

Con el comando \texttt{ip} es posible asociar una tabla de ruteo distinta con cada acceso en el escenario de Split Access. La decisión de qué tráfico utilizará cada tabla de ruteo se puede configurar en varias formas. 

\subsubsection{Creación de tablas de ruteo}
El paquete \texttt{iproute2} identifica las tablas con números, pero es más cómodo para el administrador establecer unos nombres simbólicos en el archivo \texttt{/etc/iproute2/rt\_tables}. Simplemente se eligen números y nombres arbitrarios, que no colisionen con los ya reservados. Los nombres servirán posteriormente para referirse a las diferentes tablas de ruteo usando el comando \texttt{ip}, y pueden ser los nombres de los ISP o proveedores de acceso. Una vez creados los nombres simbólicos, por cada acceso a Internet disponible es necesario identificar los siguientes datos.
\begin{itemize}
	\item Nombre de la tabla de ruteo ($P$).
	\item Dirección del router del proveedor ($R$).
	\item Dirección de red donde se ubica el router ($N$).
	\item Interfaz del router de la organización que comunica con el router del proveedor ($IF$).
	\item Dirección de dicha interfaz ($IP$).
\end{itemize}

Con estos datos, por cada acceso disponible:
\begin{enumerate}
	\item Se insertan las reglas de ruteo en su tabla correspondiente. 
\begin{lstlisting}
ip route add N dev IF src IP table P
ip route add default via R table P
\end{lstlisting}
	\item Se inserta en la tabla \texttt{main} la regla de ruteo al gateway del acceso.
	\begin{lstlisting}
ip route add N dev IF src IP
\end{lstlisting}
\end{enumerate} 

\subsection{Definición de políticas}
A continuación se establecen las reglas por las cuales cada host, cuyo tráfico se va a reenviar, utilizará una u otra tabla de ruteo. Consideraremos tres diferentes formas de diseñar y establecer las políticas. 

\subsubsection{1. Política según el origen}
La idea es simplemente separar los hosts de las redes locales en grupos, estáticamente, por cualquier criterio conveniente; y derivar el tráfico de cada grupo por rutas diferentes y fijas. 

Por cada host o red con dirección $IP$ que se quiere conducir por el acceso $P$:
\begin{lstlisting}
ip rule add from IP table P 
\end{lstlisting}

\recuadro{Las reglas se almacenan en una lista en memoria y se aplican en el orden en que están almacenadas. La primera regla cuya cláusula \texttt{from} coincide con el origen del tráfico determina la tabla que se aplicará. 

Sin embargo, cada regla introducida se inserta al principio de la lista, de forma que las últimas reglas introducidas tienen prioridad sobre las anteriores.}

En este ejemplo, se desea que la red 10.0.0.0/24 en general use la tabla T1; \textbf{pero} la mitad superior de esos hosts (los que tienen último octeto mayor o igual que 128) deben usar la T2; \textbf{excepto}, de éstos, el host 129, que usará la tabla T1. 

\begin{itemize}
	\item Al insertar las reglas: 
\begin{lstlisting}
# ip rule add from 10.0.0.0/24 table T1
# ip rule add from 10.0.0.128/25 table T2
# ip rule add from 10.0.0.129/32 table T1
\end{lstlisting}
	\item En memoria:
	\begin{lstlisting}
# ip rule show
0:	from all lookup local 
32763:	from 10.0.0.129 lookup T1
32764:	from 10.0.0.128/25 lookup T2 
32765:	from 10.0.0.0/24 lookup T1 
32766:	from all lookup main 
32767:	from all lookup default 
\end{lstlisting}
\end{itemize}

Las reglas se insertarán de lo general a lo particular, para que, al aplicarlas, las que se encuentren primero en la lista en memoria sean las más específicas. 
 
\subsubsection{2. Política según la clase de tráfico}

Utilizando \texttt{iptables} para manipular la tabla \texttt{mangle} de Netfilter, podemos crear \textbf{marcas} numéricas, que acompañan a los paquetes en su tránsito por el kernel. Estas marcas son reconocidas por el proceso de ruteo, y sirven, entre otras cosas, para determinar la tabla de ruteo a utilizar para cada paquete. Las marcas existen únicamente durante el viaje de los paquetes dentro del sistema, y al ser emitido el paquete por su interfaz de salida desaparecen. 

Marcando los paquetes podemos definir cualquier política basada en las características que puede distinguir iptables en ellos (direcciones IP origen o destino, protocolo, etc.). Cada conjunto de valores de estas características definirá una clase de tráfico.

Por ejemplo, para el tráfico TCP al port $X$, con IP origen $S$ y destino $D$ (o cualesquiera otras condiciones que sea posible seleccionar mediante iptables), marcamos los paquetes en el ingreso con la marca numérica $Z$.
\begin{lstlisting}
iptables -t mangle -A PREROUTING -p tcp --dport X -s S -d D -j MARK --set-mark Z
\end{lstlisting}

Luego se establecen las reglas que derivan el tráfico con cada marca $Z$ a la tabla de ruteo $P$ que se desee:
\begin{lstlisting}
ip rule add fwmark Z table P
\end{lstlisting}

\subsubsection{3. Política de balance de carga}
El paquete iproute2 permite establecer una ruta default con capacidad de \textit{multipath}. Con la configuración adecuada, esta ruta por defecto tomará alternativamente cada uno de los enlaces disponibles, asociando pesos a los enlaces. 

Para cada paquete que se emita, el proceso de ruteo elegirá una ruta entre las disponibles en modo \textit{round-robin}, afectando la elección con determinados pesos. Es decir, si un enlace $A$ tiene peso doble (o triple, etc.) que otro enlace $B$, entonces $A$ será elegido el doble (o triple, etc.) de veces que $B$. 

Sin embargo, una vez que se elija automáticamente una ruta a un destino, ésta quedará en la cache de rutas. Esto es necesario para que la conexión no cambie la dirección IP de origen, pero al mismo tiempo influye en la distribución de tráfico, porque las conexiones siguientes al mismo destino utilizarán la misma ruta hasta que ésta caiga de la cache de rutas o hasta que sea vaciada administrativamente (con \texttt{ip route flush cache}). La consecuencia es que no se logrará equilibrar exactamente el uso de los enlaces en la proporción especificada, sino en forma aproximada. 
  
En este ejemplo, $R1$ y $R2$ son las direcciones de los gateways de los enlaces, $IF1$ e $IF2$ las interfaces que conectan al router con ellos, y $W1$ y $W2$ son números que designan los pesos relativos de la carga que se desea hacer circular por cada enlace. Si los enlaces contratados tienen aproximadamente la misma velocidad de transmisión, pesos iguales deben utilizar los enlaces aproximadamente en la misma proporción.


\begin{lstlisting}
ip route add default scope global nexthop via R1 dev IF1 weight W1 nexthop via R2 dev IF2 weight W2
\end{lstlisting}


% 2. En round-robin
% Existe un módulo de Netfilter que emite circularmente un número módulo N para marcar paquetes. Para cada marca se agrega una regla que dirige los paquetes marcados a una tabla de ruteo determinada. Este módulo debería aplicarse a los paquetes SYN que originan una conexión. De lo contrario se descompone la conversación en los diferentes links, lo cual no sirve para conexiones outbound.
% Después:
% 
% ip rule add fwmark Z table Z
% ip route flush cache
% 

\subsection{Temas de práctica}
\subsubsection{Laboratorio virtual de Split Access}
La organización tiene dos redes locales o VLANs, LAN1 y LAN2, conectadas por un router. Se ha contratado acceso a Internet de dos proveedores diferentes. Ambos han provisto sus routers r1 y r2 con una dirección interna privada. El problema consiste en poder ofrecer a las redes locales de la organización acceso por uno u otro de los servicios, en la forma más flexible posible.

\figura[10]{splitlab}{Esquema del laboratorio virtual de Split Access}{origen.png}

En el laboratorio se han fijado las direcciones que aparecen en el diagrama de Fig. \ref{fig:splitlab}. Todas las demás deben ser configuradas con direcciones a elección. 
\begin{enumerate} 
	\item Establecer ruteo estático según la red de origen.
	\begin{itemize}
	\item Comprobar que r1 y r2 tienen acceso a Internet. 
	\item Dar direcciones a las redes locales y establecer ruta por defecto para los hosts. 
	\item En \texttt{router} reemplazar el ruteo simple por \textit{masquerading}.	
	\item Comprobar que los hosts pc1 a pc4 tienen ruta hacia las interfaces internas de r1 y r2. 
	\item Especificar en \texttt{router} un gateway default en la tabla main. 
	\item Comprobar que los hosts pc1 a pc4 tienen acceso a sitios de Internet. 
	\item Por cada proveedor, establecer en \texttt{router} una tabla de ruteo en \texttt{/etc/iproute2/rt\_tables}. Configurar cada tabla de ruteo con regla default y especificación de IP de salida. 
	\item Agregar reglas de ruteo por origen estáticas por cada red local. 
	\item Desde los hosts, hacer ping a un host público de Internet. Ver con tcpdump en las interfaces internas de r1 y de r2 que \texttt{router}  aplica la política por origen correspondiente. 
	\end{itemize}
	\item Establecer políticas por clase de tráfico.
	\begin{itemize}
		\item Modificar la configuración de las reglas de selección de tablas, aplicando con \texttt{iptables} una política que derive el tráfico de HTTP por el acceso 1 y todo otro tráfico por el acceso 2.
		\item Establecer que el tráfico ICMP debe usar un enlace, y el tráfico TCP el otro.
	\end{itemize}
	\item Incorporar excepciones.
	\begin{itemize}
		\item Establecer la excepción de que un determinado host de la organización utilizará un acceso diferente que el resto de la red a la que pertenece.	
		\item Establecer la excepción de que un determinado servidor HTTP, externo a la organización, se accederá siempre por el acceso 2. ¿Cómo se diferencia este caso del anterior? 
	\end{itemize}
	\item Eliminar la regla de ruteo default, y agregar la directiva ip con pesos para balancear carga. 
	\begin{itemize}
		\item Experimentar con transferencias de archivos, ping, etc. 
		\item Experimentar con pesos diferentes. Ensayar transferencias desde diferentes mirrors de descargas (para que las conexiones sean hacia direcciones diferentes) y observar los bytes traficados por cada interfaz del router según muestra \texttt{ifconfig}. 
		\item ¿Es posible combinar políticas por clase de tráfico con política de balance de carga? Es decir, ¿se puede establecer una política general de balance de carga, pero además asociar estáticamente una clase particular de tráfico con un enlace fijo?
	\end{itemize}
	\item En \texttt{router}, implementar scripts o funciones que conviertan el acceso a Internet en tolerante a fallos de ISP. Paso 1: pensar el problema. Paso 2 (opcional): ver Sugerencias en las Notas del Laboratorio.
\end{enumerate}
	

	

\subsubsection{Notas del Laboratorio}
\begin{itemize}
	\item Para dar acceso a Internet a las máquinas virtuales del laboratorio que funcionan como routers, es importante dar en la consola \textbf{del host} los comandos que enmascaran el tráfico de salida proveniente de esas máquinas virtuales.
\begin{lstlisting}
iptables -t nat -A POSTROUTING -s 172.16.100.0/24 -o eth0 -j MASQUERADE
iptables -A FORWARD -i eth0 -m state --state RELATED,ESTABLISHED -j ACCEPT
\end{lstlisting}
	\item Para efectuar una transferencia de un archivo sin descargarlo en el disco se puede utilizar:
\begin{lstlisting}
wget -O - http://servidor/archivo  > /dev/null
\end{lstlisting}
\item Sugerencias de análisis para el problema de Tolerancia a Fallos
	\begin{itemize}
		\item \texttt{testping}: Monitorea periódicamente el alcance mediante ping a una dirección IP testigo. Devuelve alcanzable o no alcanzable, luego de una secuencia de $n$ pings con $m$ exitosos o fallidos. Cantidades $n$ y $m$ configurables. Argumentos: dirección IP testigo, n, m. 
		\item \texttt{monitorear}: Ciclo infinito con intervalo configurable que ejecuta a \texttt{testping}. Recibe como argumentos los nombres de dos scripts o funciones, el intervalo, la dirección IP testigo, n, m. Reconoce los eventos resultantes de \texttt{testping} y actúa en consecuencia:
		\begin{itemize}
			\item Evento de dirección testigo inalcanzable $\rightarrow$ ejecuta el primer script.
			\item Evento de dirección testigo que vuelve a ser alcanzable $\rightarrow$ ejecuta el segundo script.
		\end{itemize}
		\item \texttt{toleranciaAFallos}: Lanzado en dos instancias, una para ISP A, y otra para ISP B. Ejecuta a \texttt{monitorear} dándole como argumentos los nombres de dos funciones o scripts que contienen las reglas de ruteo para implementar el acceso por uno u otro proveedor. 
	\end{itemize}
\end{itemize}





%----------- A N E X O S ---------
\newpage
\part {Anexos}
\appendix
\subsection{iptables.log}
\label{subsec:iptables.log}
\begin{lstlisting}
 Logged 539 packets on interface eth1
   From 0000:0000:1011:1213:0100:0000:0000:0000 - 3 packets to icmpv6(130)
   From 0000:0000:0000:859e:0100:0000:0000:0000 - 1 packet to icmpv6(130)
   From 0000:0023:ff53:4d42:0100:0000:0000:0000 - 1 packet to icmpv6(130)
   From 0000:0000:0000:3433:0100:0000:0000:0000 - 1 packet to icmpv6(130)
   From 0000:0000:0000:3132:0100:0000:0000:0000 - 1 packet to icmpv6(130)
   From 0000:0000:0000:7d3a:0100:0000:0000:0000 - 1 packet to icmpv6(130)
   From 0000:0000:0000:6e63:0100:0000:0000:0000 - 1 packet to icmpv6(130)
   From 0000:0000:0000:937f:0100:0000:0000:0000 - 1 packet to icmpv6(130)
   From 0000:0000:0000:0000:0000:0000:0000:0000 - 2 packets to icmpv6(130)
   From 0000:0000:0000:0000:0100:0000:0000:0000 - 40 packets to icmpv6(130)
   From 0000:0000:0000:6569:0100:0000:0000:0000 - 1 packet to icmpv6(130)
   From 000e:175f:531c:580e:0100:0000:0000:0000 - 1 packet to icmpv6(130)
   From 0011:11db:a2d4:0a00:0100:0000:0000:0000 - 2 packets to icmpv6(130)
   From 002c:799a:0694:8c26:0100:0000:0000:0000 - 1 packet to icmpv6(130)
   From 0100:0000:0600:0000:0100:0000:0000:0000 - 2 packets to icmpv6(130)
   From 0101:080a:00c6:0621:0100:0000:0000:0000 - 1 packet to icmpv6(130)
   From 0101:080a:00c6:f565:0007:994d:0000:0000 - 1 packet to icmpv6(130)
   From 0101:080a:00c7:39ad:0100:0000:0000:0000 - 1 packet to icmpv6(130)
   From 0101:080a:00c7:3e49:0100:0000:0000:0000 - 1 packet to icmpv6(130)
   From 0101:080a:00c7:5bd1:0100:0000:0000:0000 - 1 packet to icmpv6(130)
   From 0101:080a:1551:7183:0100:0000:0000:0000 - 1 packet to icmpv6(130)
   From 0101:080a:1a9a:1543:0100:0000:0000:0000 - 1 packet to icmpv6(130)
   From 0101:080a:4553:94db:0100:0000:0000:0000 - 1 packet to icmpv6(130)
   From 0101:080a:4557:126c:0100:0000:0000:0000 - 1 packet to icmpv6(130)
   From 0101:080a:4559:6ffd:0100:0000:0000:0000 - 1 packet to icmpv6(130)
   From 0101:080a:4559:e9ac:0100:0000:0000:0000 - 1 packet to icmpv6(130)
   From 0101:080a:455a:3fd8:0100:0000:0000:0000 - 1 packet to icmpv6(130)
   From 0101:080a:455a:cc78:0100:0000:0000:0000 - 1 packet to icmpv6(130)
   From 0101:080a:455c:c658:0100:0000:0000:0000 - 1 packet to icmpv6(130)
   From 0101:080a:455d:4147:0100:0000:0000:0000 - 1 packet to icmpv6(130)
   From 0101:080a:455d:bbfc:0100:0000:0000:0000 - 1 packet to icmpv6(130)
   From 0101:080a:455e:3351:0100:0000:0000:0000 - 1 packet to icmpv6(130)
   From 0101:080a:4567:104c:0100:0000:0000:0000 - 1 packet to icmpv6(130)
   From 0101:080a:4575:8fa3:0100:0000:0000:0000 - 1 packet to icmpv6(130)
   From 0101:080a:4575:fcd3:0100:0000:0000:0000 - 1 packet to icmpv6(130)
   From 0101:080a:4584:d388:0100:0000:0000:0000 - 1 packet to icmpv6(130)
   From 0101:080a:458c:f368:0100:0000:0000:0000 - 1 packet to icmpv6(130)
   From 0101:080a:4594:8809:0100:0000:0000:0000 - 1 packet to icmpv6(130)
   From 0102:0417:0000:0000:0100:0000:0000:0000 - 1 packet to icmpv6(130)
   From 0102:a16d:0a00:00c8:0100:0000:0000:0000 - 1 packet to icmpv6(130)
   From 0204:05b4:0402:080a:0100:0000:0000:0000 - 2 packets to icmpv6(130)
   From 0300:0000:0400:0000:0100:0000:0000:0000 - 3 packets to icmpv6(130)
   From 036b:696d:0675:6e63:0100:0000:0000:0000 - 1 packet to icmpv6(130)
   From 10d4:d5cb:f80c:14d5:0100:0000:0000:0000 - 1 packet to icmpv6(130)
   From 1400:0300:9709:0000:0100:0000:0000:0000 - 1 packet to icmpv6(130)
   From 2269:6422:3a20:2232:0100:0000:0000:0000 - 1 packet to icmpv6(130)
   From 3037:3337:3431:3832:0100:0000:0000:0000 - 1 packet to icmpv6(130)
   From 3135:3332:3a38:3439:0100:0000:0000:0000 - 1 packet to icmpv6(130)
   From 3234:3238:323a:3834:0100:0000:0000:0000 - 1 packet to icmpv6(130)
   From 3333:3238:323a:3834:0100:0000:0000:0000 - 1 packet to icmpv6(130)
   From 3335:3839:323a:3834:0100:0000:0000:0000 - 1 packet to icmpv6(130)
   From 3336:3532:323a:3834:0100:0000:0000:0000 - 1 packet to icmpv6(130)
   From 3336:3837:3039:3132:0100:0000:0000:0000 - 8 packets to icmpv6(130)
   From 3430:3734:3330:3532:0100:0000:0000:0000 - 1 packet to icmpv6(130)
   From 3432:3230:3239:3a33:0100:0000:0000:0000 - 1 packet to icmpv6(130)
   From 3432:3635:3239:3a33:0100:0000:0000:0000 - 1 packet to icmpv6(130)
   From 3433:3230:3332:3a38:0100:0000:0000:0000 - 1 packet to icmpv6(130)
   From 3433:3534:3039:3a33:0100:0000:0000:0000 - 1 packet to icmpv6(130)
   From 3734:3237:3339:313a:0100:0000:0000:0000 - 1 packet to icmpv6(130)
   From 3734:3330:3238:313a:0100:0000:0000:0000 - 1 packet to icmpv6(130)
   From 3734:3330:3636:313a:0100:0000:0000:0000 - 1 packet to icmpv6(130)
   From 3734:3330:3930:323a:0100:0000:0000:0000 - 1 packet to icmpv6(130)
   From 3734:3334:3533:313a:0100:0000:0000:0000 - 1 packet to icmpv6(130)
   From 3734:3334:3736:393a:0100:0000:0000:0000 - 1 packet to icmpv6(130)
   From 616a:6f72:223a:2031:0100:0000:0000:0000 - 1 packet to icmpv6(130)
   From 6d65:6e74:7322:3a7b:0100:0000:0000:0000 - 9 packets to icmpv6(130)
   From 6f20:3134:3037:3433:0100:0000:0000:0000 - 2 packets to icmpv6(130)
   From 7473:223a:7b7d:7d00:0100:0000:0000:0000 - 1 packet to icmpv6(130)
   From 7473:223a:7b7d:7d32:0100:0000:0000:0000 - 1 packet to icmpv6(130)
   From 7473:223a:7b7d:7d3a:0100:0000:0000:0000 - 4 packets to icmpv6(130)
   From 7473:223a:7b7d:7d7b:0100:0000:0000:0000 - 3 packets to icmpv6(130)
   From 756e:6e69:6e67:223a:0100:0000:0000:0000 - 1 packet to icmpv6(130)
   From bf1d:f661:984e:fcb0:0100:0000:0000:0000 - 1 packet to icmpv6(130)
   From c011:fdda:43fe:4d59:65fe:e1aa:c7d5:683a - 1 packet to icmpv6(130)
   From c012:aa5c:173c:261c:0100:0000:0000:0000 - 1 packet to icmpv6(130)
   From c012:cfa3:e641:3b5f:0100:0000:0000:0000 - 1 packet to icmpv6(130)
   From c013:4b15:1c15:3311:0100:0000:0000:0000 - 1 packet to icmpv6(130)
   From c013:9389:bcf8:f7d4:0100:0000:0000:0000 - 1 packet to icmpv6(130)
   From c014:ff7d:0000:0000:0100:0000:0000:0000 - 1 packet to icmpv6(130)
   From e063:837f:0000:0000:0100:0000:0000:0000 - 3 packets to icmpv6(130)
   From e063:837f:2077:937f:0100:0000:0000:0000 - 1 packet to icmpv6(130)
   From e063:837f:301f:937f:0100:0000:0000:0000 - 3 packets to icmpv6(130)
   From 0.0.0.0 - 154 packets to igmp(0)
   From 10.0.3.4 - 154 packets to igmp(0)
   From 10.0.3.21 - 75 packets to igmp(0)
   From 10.0.3.209 - 2 packets to igmp(0)

 Listed by source hosts:
 Logged 18 packets on interface virbr0
   From fe80:0000:0000:0000:5054:00ff:feed:8246 - 18 packets to udp(5353)

 Listed by source hosts:
 Logged 50 packets on interface wlan0
   From 10.0.4.1 - 2 packets to udp(68)
   From 64.233.186.188 - 15 packets to tcp(44421,53418)
   From 64.235.151.8 - 5 packets to tcp(56214)
   From 173.194.42.0 - 1 packet to tcp(56677)
   From 173.194.42.21 - 3 packets to tcp(58597)
   From 173.194.42.22 - 7 packets to tcp(35536)
   From 173.194.42.75 - 3 packets to tcp(48585)
   From 173.194.42.85 - 4 packets to tcp(44550)
   From 173.194.42.86 - 1 packet to tcp(51940)
   From 192.168.1.1 - 2 packets to udp(68)
   From 192.168.2.1 - 2 packets to udp(68)
   From 195.135.221.134 - 1 packet to tcp(51756)
   From 195.154.174.66 - 1 packet to tcp(58047)
   From 200.42.136.212 - 3 packets to tcp(59351,59361)
\end{lstlisting}

%------------------------------------------------------------

\end{document}
